\chaptermark{Prestack migration}

\chapter{Prestack migration}

We have formulated basic migration concepts in terms of zero-offset data, even though real data are anything but zero-offset. Fortunately, the ideas are generalizable to finite-offset data.

\section{Prestack migration concepts}

Kirchhoff migration\index{migration!Kirchhoff} is perfectly amenable to prestack data\index{migration!prestack}. As long as one can calculate (or tabulate) the traveltimes\index{traveltime!two way} from source to reflection
point to receiver, one can sum into a point all the data that might have come from that point, or, alternatively, broadcast a data point into all
the locations from which the data might have come. Kirchhoff methods are happy with restricted data sets such as constant offset, though in some
circumstances they may produce artifacts. There are of course issues of amplitude, phase, aliasing, etc., which must be dealt with, some of which are discussed below.

\begin{examples}
A prestack difference-equation migration algorithm\index{migration!prestack difference equation} requires extrapolation of both sources and receivers. Given full surface coverage, one could
do this by applying the wave equation separately to source and receiver coordinates, downward continuing each in turn, and taking the time-zero doubly
downward-continued data to be the migrated image. To image a restricted data set, one must be more devious. For example, a constant-source
set can be imaged by downward-continuing the receiver coordinates and cross-correlating at each image point with a forward extrapolation of the source
impulse.
\end{examples}

The simplicity of downward continuation in \textit{f--k} space makes wavenumber--frequency migration methods particularly easy to construct and explain.
However, these methods are not amenable to migration of constant-offset or most other partial data sets. Fourier transforms over both source
and receiver locations are required, which implies that for every source, there must be a complete set of receivers, and vice versa. Fourier transforms over both source
and receiver locations are required, which implies that for every source, there must be a complete set of receivers, and vice versa. By invoking
asymptotic approximations, the demand for completeness can be relaxed somewhat, but the full flexibility enjoyed by other methods is elusive. By invoking
asymptotic approximations, the demand for completeness can be relaxed somewhat, but the full flexibility enjoyed by other methods is elusive.

\begin{examples}
A prestack difference-equation migration algorithm\index{migration!prestack difference equation} requires extrapolation of both sources and receivers. Given full surface coverage, one could
do this by applying the wave equation separately to source and receiver coordinates, downward continuing each in turn, and taking the time-zero doubly
downward-continued data to be the migrated image. To image a restricted data set, one must be more devious. For example, a constant-source
set can be imaged by downward-continuing the receiver coordinates and cross-correlating at each image point with a forward extrapolation of the source
impulse.
\end{examples}

The simplicity of downward continuation in \textit{f--k} space makes wavenumber--frequency migration methods particularly easy to construct and explain.
However, these methods are not amenable to migration of constant-offset or most other partial data sets. Fourier transforms over both source
and receiver locations are required, which implies that for every source, there must be a complete set of receivers, and vice versa. By invoking
asymptotic approximations, the demand for completeness can be relaxed somewhat, but the full flexibility enjoyed by other methods is elusive. By invoking
asymptotic approximations, the demand for completeness can be relaxed somewhat, but the full flexibility enjoyed by other methods is elusive.

\begin{problems}
A prestack difference-equation migration algorithm\index{migration!prestack difference equation} requires extrapolation of both sources and receivers. Given full surface coverage, one could
do this by applying the wave equation separately to source and receiver coordinates, downward continuing each in turn, and taking the time-zero doubly
downward-continued data to be the migrated image. To image a restricted data set, one must be more devious. For example, a constant-source
set can be imaged by downward-continuing the receiver coordinates and cross-correlating at each image point with a forward extrapolation of the source
impulse.
\end{problems}

The simplicity of downward continuation in \textit{f--k} space makes wavenumber--frequency migration methods particularly easy to construct and explain.
However, these methods are not amenable to migration of constant-offset or most other partial data sets. Fourier transforms over both source
and receiver locations are required, which implies that for every source, there must be a complete set of receivers, and vice versa. By invoking
asymptotic approximations, the demand for completeness can be relaxed somewhat, but the full flexibility enjoyed by other methods is elusive.

\begin{problems}
A prestack difference-equation migration algorithm\index{migration!prestack difference equation} requires extrapolation of both sources and receivers. Given full surface coverage, one could
do this by applying the wave equation separately to source and receiver coordinates, downward continuing each in turn, and taking the time-zero doubly
downward-continued data to be the migrated image. To image a restricted data set, one must be more devious. For example, a constant-source
set can be imaged by downward-continuing the receiver coordinates and cross-correlating at each image point with a forward extrapolation of the source
impulse.
\end{problems}

The simplicity of downward continuation in \textit{f--k} space makes wavenumber--frequency migration methods particularly easy to construct and explain.
However, these methods are not amenable to migration of constant-offset or most other partial data sets. Fourier transforms over both source
and receiver locations are required, which implies that for every source, there must be a complete set of receivers, and vice versa. By invoking
asymptotic approximations, the demand for completeness can be relaxed somewhat, but the full flexibility enjoyed by other methods is elusive. Fourier transforms over both source
and receiver locations are required, which implies that for every source, there must be a complete set of receivers, and vice versa. By invoking
asymptotic approximations, the demand for completeness can be relaxed somewhat, but the full flexibility enjoyed by other methods is elusive.

\begin{framed}
\noindent The simplicity of downward continuation in \textit{f--k} space makes wavenumber--frequency migration methods particularly easy to construct and explain.
However, these methods are not amenable to migration of constant-offset or most other partial data sets. Fourier transforms over both source
and receiver locations are required, which implies that for every source, there must be a complete set of receivers, and vice versa. By invoking
asymptotic approximations, the demand for completeness can be relaxed somewhat, but the full flexibility enjoyed by other methods is elusive.

The simplicity of downward continuation in \textit{f--k} space makes wavenumber--frequency migration methods particularly easy to construct and explain.
However, these methods are not amenable to migration of constant-offset or most other partial data sets. Fourier transforms over both source
and receiver locations are required, which implies that for every source, there must be a complete set of receivers, and vice versa. By invoking
asymptotic approximations, the demand for completeness can be relaxed somewhat, but the full flexibility enjoyed by other methods is elusive.

The simplicity of downward continuation in \textit{f--k} space makes wavenumber--frequency migration methods particularly easy to construct and explain.
However, these methods are not amenable to migration of constant-offset or most other partial data sets. Fourier transforms over both source
and receiver locations are required, which implies that for every source, there must be a complete set of receivers, and vice versa. By invoking
asymptotic approximations, the demand for completeness can be relaxed somewhat, but the full flexibility enjoyed by other methods is elusive.

The simplicity of downward continuation in \textit{f--k} space makes wavenumber--frequency migration methods particularly easy to construct and explain.
However, these methods are not amenable to migration of constant-offset or most other partial data sets. Fourier transforms over both source
and receiver locations are required, which implies that for every source, there must be a complete set of receivers, and vice versa. By invoking
asymptotic approximations, the demand for completeness can be relaxed somewhat, but the full flexibility enjoyed by other methods is elusive.
\end{framed}

The simplicity of downward continuation in \textit{f--k} space makes wavenumber--frequency migration methods particularly easy to construct and explain.
However, these methods are not amenable to migration of constant-offset or most other partial data sets. Fourier transforms over both source
and receiver locations are required, which implies that for every source, there must be a complete set of receivers, and vice versa. By invoking
asymptotic approximations, the demand for completeness can be relaxed somewhat, but the full flexibility enjoyed by other methods is elusive.

\begin{theorem}
The simplicity of downward continuation in \textit{f--k} space makes wavenumber--frequency migration methods particularly easy to construct and explain.
However, these methods are not amenable to migration of constant-offset or most other partial data sets. Fourier transforms over both source
and receiver locations are required, which implies that for every source, there must be a complete set of receivers, and vice versa. By invoking
asymptotic approximations, the demand for completeness can be relaxed somewhat, but the full flexibility enjoyed by other methods is elusive.
\end{theorem}

The simplicity of downward continuation in \textit{f--k} space makes wavenumber--frequency migration methods particularly easy to construct and explain.
However, these methods are not amenable to migration of constant-offset or most other partial data sets. Fourier transforms over both source
and receiver locations are required, which implies that for every source, there must be a complete set of receivers, and vice versa. By invoking
asymptotic approximations, the demand for completeness can be relaxed somewhat, but the full flexibility enjoyed by other methods is elusive.

\begin{definition}
The simplicity of downward continuation in \textit{f--k} space makes wavenumber--frequency migration methods particularly easy to construct and explain.
However, these methods are not amenable to migration of constant-offset or most other partial data sets. Fourier transforms over both source
and receiver locations are required, which implies that for every source, there must be a complete set of receivers, and vice versa. By invoking
asymptotic approximations, the demand for completeness can be relaxed somewhat, but the full flexibility enjoyed by other methods is elusive.
\end{definition}

\section{Prestack \textit{f--k} migration in two dimensions}\index{migration!2-D prestack $f{-}k$}

The simplicity of downward continuation in \textit{f--k} space makes wavenumber--frequency migration methods particularly easy to construct and explain.
However, these methods are not amenable to migration of constant-offset or most other partial data sets. Fourier transforms over both source
and receiver locations are required, which implies that for every source, there must be a complete set of receivers, and vice versa. By invoking
asymptotic approximations, the demand for completeness can be relaxed somewhat, but the full flexibility enjoyed by other methods is elusive.

Given a full 2-D data set \(D\left(x_g|x_s;t\right)\), where \(x_s\) is source location, \(x_g\) is receiver location, and \textit{t} is time, one
can perform a Fourier transform over all coordinates:
\begin{framed}
\begin{equation}\text{\textit{$D$}}\left(\text{\textit{$k$}}_{\text{\textit{gx}}}|\text{\textit{$k$}}_{\text{\textit{sx}}};\omega \right) =\text{
 }\int d\text{\textit{$x$}}_{\text{\textit{$g$}}}\int d\text{\textit{$x$}}_{\text{\textit{$s$}}}\int d\text{\textit{$t$}}\text{\textit{$ $}}\text{\textit{$D$}}\left(\text{\textit{$x$}}_{\text{\textit{$g$}}}|\text{\textit{$x$}}_{\text{\textit{$s$}}};\text{\textit{$t$}}\right)e^{i
\left(\text{\textit{$k$}}_{\text{\textit{sx}}}\text{\textit{$x$}}_{\text{\textit{$s$}}}-\text{\textit{$k$}}_{\text{\textit{gx}}}\text{\textit{$x$}}_{\text{\textit{$g$}}}+\omega
 \text{\textit{$t$}}\right)}.\end{equation}
\end{framed}

In the convention employed here, the forward time transform is performed with positive phase, as is the forward source-coordinate transform. The
receiver-coordinate transform is performed with negative phase. For a discussion of these conventions, see Appendix C.

\subsubsection{Downward continuation}

To perform downward continuation, we consider the data to be the expression of a wave recorded at \(\left(x_g,0\right)\) produced by a source at
\(\left(x_s,0\right)\), and reflected from multiple subsurface locations (\textit{x}, \textit{z}). Denoting this wave as \(P\left(x_g,z_g|x_s,z_s;\omega
\right)\), we have
\begin{equation}
\text{\textit{$D$}}\left(\text{\textit{$x$}}_{\text{\textit{$g$}}}|\text{\textit{$x$}}_{\text{\textit{$s$}}};\text{\textit{$t$}}\right)=\text{\textit{$P$}}\left(\text{\textit{$x_g$}},0\left|\text{\textit{$x_s$}},\right.0;\text{\textit{$t$}}\right),\end{equation}
and, in the frequency--wavenumber domain,
\begin{equation}\text{\textit{$D$}}\left(\text{\textit{$k$}}_{\text{\textit{gx}}}|\text{\textit{$k$}}_{\text{\textit{sx}}};\omega \right)=\text{\textit{$P$}}\left(\text{\textit{$k$}}_{\text{\textit{gx}}},0\left|\text{\textit{$k$}}_{\text{\textit{sx}}},\right.0;\omega
\right).\end{equation}

(While we have written these equations for 2-D data, the 3-D expression is a simple extension.) Analogous to Section 2.6,
the Fourier transforms isolate a single plane-wave component of both the wave from the source to the reflectors and the wave from the reflectors
to the receiver. Assuming a constant velocity \textit{c}, the wavenumber vector\index{wavenumber!vector} of the (downgoing) source wave\index{wave!downgoing} is
\begin{equation}\tilde{\boldsymbol{k}}_s =\left(\text{\textit{$k_{\text{sx}}$}},\text{\textit{$k_{\text{sz}}$}}\right),\end{equation}
where
\begin{equation}\text{\textit{$k$}}_{\text{\textit{sz}}}=+\frac{\omega }{\text{\textit{$c$}}}\text{  }\sqrt{1-\frac{\text{\textit{$k$}}_{\text{\textit{sx}}}{}^2\text{\textit{$c$}}^2}{\omega
^2}},\end{equation}
and that of the (upgoing) receiver wave\index{wave!upgoing} is
\begin{equation}\tilde{\boldsymbol{k}}_g =\left(\text{\textit{$k_{\text{gx}}$}},\text{\textit{$k_{\text{gz}}$}}\right),\end{equation}
with
\begin{equation}\text{\textit{$k$}}_{\text{\textit{gz}}}=- \frac{\omega }{\text{\textit{$c$}}} \sqrt{1-\frac{\text{\textit{$k$}}_{\text{\textit{gx}}}{}^2\text{\textit{$c$}}^2}{\omega
^2}}.\end{equation}

As for poststack migration, we confine attention to the physically propagating wavenumbers\index{wavenumber}, where the arguments of both square roots are positive and the vertical wavenumbers are real.

To downward-continue both the source and receiver waves to the same depth \textit{z}, we simply multiply the data by the phase term \(e^{i \left(k_{\text{\textit{gz}}}-k_{\text{\textit{sz}}}\right)z}\):
\begin{equation}\text{\textit{$P$}}\left(\text{\textit{$k$}}_{\text{\textit{gx}}},\text{\textit{$z$}}\left|\text{\textit{$k$}}_{\text{\textit{sx}}},\right.\text{\textit{$z$}};\omega
\right) =\text{  }\text{\textit{$D$}}\left(\text{\textit{$k$}}_{\text{\textit{gx}}}|\text{\textit{$k$}}_{\text{\textit{sx}}};\omega \right) e^{i\left(\text{\textit{$k$}}_{\text{\textit{gz}}}-\text{\textit{$k$}}_{\text{\textit{sz}}}\right)\text{\textit{$z$}}}.\end{equation}

Following the convention of Appendix C for Fourier transforms, the source\index{wavenumber!source} and receiver\index{wavenumber!receiver} wavenumber phases are given opposite signs. This is because
downward continuation extrapolates the source wave forward and the receiver wave backwards.

Defining \(k_z\) to be the total vertical wavenumber\index{wavenumber!vertical} for the downward-continued data, we have
\begin{equation}\text{\textit{$k$}}_{\text{\textit{$z$}}}\equiv \text{\textit{$k$}}_{\text{\textit{gz}}}-\text{\textit{$k$}}_{\text{\textit{sz}}}\text{
 }=- \frac{\omega }{\text{\textit{$c$}}} \left( \sqrt{1-\frac{\text{\textit{$k$}}_{\text{\textit{gx}}}{}^2\text{\textit{$c$}}^2}{\omega ^2}} +\sqrt{1-\frac{\text{\textit{$k$}}_{\text{\textit{sx}}}{}^2\text{\textit{$c$}}^2}{\omega
^2}}\right).\end{equation}

Equation (3.9) is, for obvious reasons, often called the \textit{double-square-root} equation\index{double square root equation}. As shown in Figure 3.1, vertical wavenumber\index{wavenumber!vertical}
\(k_z\) is real within the \textit{propagating region} \(-\omega /c<k_{\text{\textit{gx}}}<\omega /c\), \(-\omega /c<k_{\text{\textit{sx}}}<\omega
/c\). For small horizontal wavenumbers\index{wavenumber!horizontal} its value is near \(-2\omega /c\). When both \(\left|k_{\textit{sx}}c/\omega \right|\) and \(\left|k_{\textit{gx}}c/\omega
\right|\) approach 1, \(k_z\) drops to zero.  When both horizontal wavenumbers exceed this limit, \(k_z\) becomes imaginary. If one wavenumber does, \(k_z\) becomes complex.

\subsubsection{Imaging}

Analogous to Section 2.6, we would like to define migration in terms of the time-zero component of the downward-continued data. The extra dimension
in prestack data complicates things a little. We must set the downward-continued source and receiver locations to be equal, so that the upgoing
and downgoing waves coincide. We would also like a migration operator that does not change the phase or amplitude spectrum of flat reflectors.
In two dimensions, there is some ambiguity with this goal. For 2-D point (more precisely, line) sources and receivers in a 2-D
Earth, the impulse response is a Hankel function\index{Hankel function} (Appendix E, Section E.10), not a delta function. Hence, if the image function for a flat reflector
is zero phase, the \nobreak corresponding data function will be the $45^{\circ}$ Hankel function. On the other hand, for real-Earth data, even where
the Earth may be approximated as 2-D, the sources and receivers should be considered points in three dimensions. In this 2.5-dimensional setting, the impulse response is a delta function, hence for flat reflectors we would like a zero-phase image function to correspond
to a zero-phase data function. To achieve the desired phase, we define the image function \textit{R}(\textit{x},\textit{z}) to be proportional
to a half-integral filtered\index{filter!half-integral} (Appendix E, Section E.6) downward-continued wave at \(x_s=x_g\equiv x,\) \(z_g=z_s=z\), and \textit{t} = 0 (note that source
and receiver must converge on the same point):
\begin{eqnarray}
\text{\textit{$R$}}(\text{\textit{$x,z$}}) \equiv \left(\Theta _{\text{\textit{$h$}}+}\text{\textit{$*$}} \text{\textit{$P$}}\text{\textit{$)$}}\right.\left(\text{\textit{$x_g$}}
=\text{\textit{$x$}},\text{\textit{$z_g$}}=\text{\textit{$z$}}\left|\text{\textit{$x_s$}}\right.=
\text{\textit{$x$}},\text{\textit{$z_s$}}=\text{\textit{$z$}};\text{\textit{$t$}}=0\right)
\nonumber\\
=\frac{1}{(2\pi )^3} \int \frac{d\omega }{\sqrt{-i \omega }} \int d\text{\textit{$k$}}_{\text{\textit{sx}}} \int d\text{\textit{$k$}}_{\text{\textit{gx}}}\text{\textit{$D$}}\left(\text{\textit{$k$}}_{\text{\textit{gx}}}|\text{\textit{$k$}}_{\text{\textit{sx}}};\omega
\right) e^{i\left(\text{\textit{$k$}}_{\text{\textit{$z$}}}\text{\textit{$z$}}+\text{\textit{$k$}}_{\text{\textit{$x$}}}\text{\textit{$x$}}\right)},
\end{eqnarray}
with \(k_z\) as defined in equation (3.9), and
\begin{equation}\text{\textit{$k$}}_{\text{\textit{$x$}}}\equiv \text{\textit{$k$}}_{\text{\textit{gx}}}-\text{\textit{$k$}}_{\text{\textit{sx}}}
.\end{equation}
with \(k_h\) the lateral offset wavenumber\index{wavenumber!offset}
\begin{equation}\text{\textit{$k$}}_{\text{\textit{$h$}}}\equiv \text{\textit{$k$}}_{\text{\textit{gx}}}+\text{\textit{$k$}}_{\text{\textit{sx}}}
.\end{equation}

Two of the three integrals in equation (3.10) can be converted to inverse Fourier transforms with a change of integration variable. We choose
as a new set of \nobreak integration variables, lateral offset wavenumber\index{wavenumber!offset} \(k_h\) and the image-point wavenumber vector\index{wavenumber!image point}
\begin{equation}
\tilde{\boldsymbol{k}}\equiv \left(\text{\textit{$k$}}_{\text{\textit{$x$}}},\text{\textit{$k$}}_{\text{\textit{$z$}}}\right).\end{equation}

The Jacobian of this change of variables is
\begin{equation}\left|\text{Det}\left[\frac{\partial \left(\text{\textit{$k_x$}},\text{\textit{$k_h$}},\text{\textit{$k_z$}}\right)}{\partial \left(\text{\textit{$k_{\text{sx}}$}},\text{\textit{$k_{\text{gx}}$}},\omega
\right)}\right]\right|=\frac{2 \omega  \text{\textit{$k_z$}}}{\text{\textit{$c$}}^2\text{\textit{$k_{\text{gz}}$}}\text{\textit{$k_{\text{sz}}$}}}.\end{equation}

Thus,
\begin{equation}\text{\textit{$R$}}(\text{\textit{$x,z$}}) = \frac{\text{\textit{$c$}}^2}{2(2\pi )^3}\text{  }\int d\text{\textit{$k$}}_{\text{\textit{$z$}}}\int
d\text{\textit{$k$}}_{\text{\textit{$x$}}}e^{i(\text{\textit{$k_xx$}}+\text{\textit{$k_zz$}})}\int d\text{\textit{$k_h$}}\frac{ \text{\textit{$k_{\text{gz}}k_{\text{sz}}$}}}{\sqrt{-i
\omega } \omega  \text{\textit{$k_z$}}}\text{\textit{$D$}}\left(\text{\textit{$k$}}_{\text{\textit{gx}}}|\text{\textit{$k$}}_{\text{\textit{sx}}};\omega
\right) .\end{equation}

The outer two integrals are now inverse Fourier transforms from \(k_z\) to \textit{z} and \(k_x\) to \(x\). Consequently, the double Fourier
transform of \textit{R} is proportional to the remaining integral over \(k_h\):
\begin{equation}\text{\textit{$R$}}\left(\text{\textit{$k_x$}},\text{\textit{$k_z$}}\right) =\frac{\text{\textit{$c$}}^2}{4\pi } \int d\text{\textit{$k$}}_{\text{\textit{$h$}}}\text{\textit{$D$}}\left(\text{\textit{$k$}}_{\text{\textit{gx}}}|\text{\textit{$k$}}_{\text{\textit{sx}}};\omega
\right) \frac{\text{\textit{$k$}}_{\text{\textit{gz}}}\text{\textit{$k$}}_{\text{\textit{sz}}}}{\sqrt{-i \omega } \omega  \text{\textit{$k_z$}}}.\end{equation}


Expression (3.16) is an integral over offset wavenumber \(k_h\). Within the integral, \(k_x\) and \(k_z\) are fixed, while \(k_{\text{\textit{gx}}}\),
\(k_{\text{\textit{sx}}}\), and $\omega $ vary with \(k_h\). The relations between the variables can be expressed as follows. Source and receiver
wavenumbers are \(k_{\text{\textit{gx}}}=\left.\left(k_x+k_h\right)\right/2\), and \(k_{\text{\textit{sx}}}=\left.\left(k_h-k_x\right)\right/2\).
Within the propagating region, the double-square-root equation\index{double square root equation} (3.9) can be solved for $\omega $ as a function of \(k_z\), \(k_x\), and \(k_h\)
(see Figure 3.1): Expression (3.16) is an integral over offset wavenumber \(k_h\). Within the integral, \(k_x\) and \(k_z\) are fixed, while \(k_{\text{\textit{gx}}}\),
\(k_{\text{\textit{sx}}}\), and $\omega $ vary with \(k_h\). The relations between the variables can be expressed as follows. Source and receiver
wavenumbers are \(k_{\text{\textit{gx}}}=\left.\left(k_x+k_h\right)\right/2\), and \(k_{\text{\textit{sx}}}=\left.\left(k_h-k_x\right)\right/2\).
\begin{equation}\omega = -\frac{\text{\textit{$c$}}\text{\textit{$ $}}\text{\textit{$k$}}_{\text{\textit{$z$}}}}{2}\sqrt{\left(1+\text{\textit{$k$}}_{\text{\textit{$x$}}}{}^2/\text{\textit{$k$}}_z{}^2\right)\left(1+\text{\textit{$k$}}_{\text{\textit{$h$}}}{}^2/\text{\textit{$k$}}_{\text{\textit{$z$}}}{}^2\right)}
.\end{equation}

The source and receiver vertical wavenumbers\index{wavenumber!vertical} \(k_{\text{\textit{sz}}}\) and \(k_{\text{\textit{gz}}}\) can also be expressed in terms of \(k_z\),
\(k_x\), and \(k_h\):
\begin{equation}\text{\textit{$k$}}_{\text{\textit{sz}}}=-\frac{\text{\textit{$k$}}_{\text{\textit{$z$}}}}{2} \left(1-\frac{\text{\textit{$k$}}_{\text{\textit{$x$}}}\text{\textit{$k$}}_{\text{\textit{$h$}}}}{\text{\textit{$k$}}_{\text{\textit{$z$}}}{}^2}\right),\end{equation}
and
\begin{equation}\text{\textit{$k$}}_{\text{\textit{gz}}}=\frac{\text{\textit{$k$}}_{\text{\textit{$z$}}}}{2} \left(1+\frac{\text{\textit{$k$}}_{\text{\textit{$x$}}}\text{\textit{$k$}}_{\text{\textit{$h$}}}}{\text{\textit{$k$}}_{\text{\textit{$z$}}}{}^2}\right).\end{equation}

\subsection{Constant-angle \textit{f--k} migration}\index{migration!constant angle $f{-}k$}

The expression (3.16) has no pretensions of true amplitude. In fact, it has one obvious deficiency. Depending on which physical properties
are changing at a point, a realizable reflection coefficient will depend on angle of incidence $\gamma $, whereas the quantity \textit{R} in (3.16)
has no angle dependence. We suspect that the angle dependence has been smothered by the integration over offset wavenumber. This integration
occurred because we wanted sources and receivers to be downward-continued to the same point, so that offset at the reflection point must be zero.
However, something subtle was missed in the argument above: offset must go to zero \textit{in the limit} as time approaches zero, and limits
do funny things sometimes. For example, a unit step function $\Theta $(\textit{x}) has the limit zero as \textit{x} approaches zero from below,
but the limit one if \textit{x} approaches zero from above. Similarly, reflectivity,\index{reflectivity} defined as the limit of the downward-continued data as time
and offset approach zero, will depend on the \textit{direction} from which the downward-continued waves approach zero (Clayton and Stolt, 1981).

We accommodate the directional dependence as follows. Changing integration variable from \(k_h\) to angle of incidence $\gamma $, equation (3.16) becomes
\begin{equation}\text{\textit{$R$}}\left(\text{\textit{$k_x$}},\text{\textit{$k_z$}}\right) =\frac{\text{\textit{$c$}}^2}{4\pi } \int d\gamma \text{\textit{$D$}}\left(\text{\textit{$k$}}_{\text{\textit{gx}}}|\text{\textit{$k$}}_{\text{\textit{sx}}};\omega
\right)\frac{\text{\textit{$k$}}_{\textit{gz}}\text{\textit{$k$}}_{\textit{sz}}}{\sqrt{-i \omega } \omega  \text{\textit{$k$}}_z}\left|\frac{d\text{\textit{$k_h$}}}{d\gamma
}\right|.\end{equation}

For an infinite-aperture migration, $\gamma$ would range from $-\pi $/2 to $\pi $/2. We can confine attention to a specific angle of incidence $\gamma$, by writing (3.20) as
\begin{equation}\text{\textit{$R$}}\left(\text{\textit{$k_x$}},\text{\textit{$k_z$}}\right) \equiv  \int d\gamma \text{\textit{$\mathit{r}$}}\left(\text{\textit{$k_x$}},\text{\textit{$k_z$}},\gamma
\right),\end{equation}
where
\begin{equation}\text{\textit{$\mathit{r}$}}\left(\text{\textit{$k_x$}},\text{\textit{$k_z$}}\text{\textit{$,$}}\gamma \right) =\frac{\text{\textit{$c$}}^2}{4\pi
} \text{\textit{$D$}}\left(\text{\textit{$k$}}_{\text{\textit{gx}}}|\text{\textit{$k$}}_{\text{\textit{sx}}};\omega \right) \frac{\text{\textit{$k$}}_{\textit{gz}}\text{\textit{$k$}}_{\textit{sz}}}{\sqrt{-i
\omega } \omega  \text{\textit{$k$}}_z}\left|\frac{d\text{\textit{$k_h$}}}{d\gamma }\right|.\end{equation}

We have assumed a relation exists for incident angle\index{angle!incident} in terms of the known parameters. To find one, define propagation angles\index{angle!propagation} \(\phi _i\) and \(\phi
_r\) for the source and receiver waves relative to the vertical as
\begin{equation}
\tilde{\boldsymbol{k}}_s =\left(\text{\textit{$k_{\text{sx}}$}},\text{\textit{$k_{\text{sz}}$}}\right)=
\frac{\omega }{\text{\textit{$c$}}}\left(-\sin  \phi _{\text{\textit{$i$}}},\cos  \phi _{\text{\textit{$i$}}}\right),
\end{equation}
\begin{equation}
\tilde{\boldsymbol{k}}_{\text{\textit{$g$}}}=\left(\text{\textit{$k_{\text{gx}}$}},\text{\textit{$k_{\text{gz}}$}}\right)
=- \frac{\omega }{\text{\textit{$c$}}}\left(\sin  \phi _{\text{\textit{$r$}}},\cos  \phi _{\text{\textit{$r$}}}\right),
\end{equation}

As defined, these two angles are positive for travel from right to left, as illustrated in Figure 3.2. For specular reflectors\index{reflections!specular}, these two angles
can be related to dip angle\index{angle!dip} $\alpha $ and angle of incidence\index{angle!incident} $\gamma$ as depicted in Figure 3.2:

From the definitions (3.9), (3.23) and (3.24), it follows that
\begin{equation}\text{\textit{$k$}}_{\text{\textit{$z$}}}=-\frac{\omega }{\text{\textit{$c$}}} \left(\cos  \phi _{\text{\textit{$i$}}}+\cos  \phi
_{\text{\textit{$r$}}}\right).\end{equation}

With the relations (3.25) and (3.26), \(k_z\) can be expressed in terms of incident and dip angles:
\begin{equation}\text{\textit{$k$}}_{\text{\textit{$z$}}}=-\frac{2\omega }{\text{\textit{$c$}}} \cos  \alpha  \cos  \gamma .\end{equation}

We also have from (3.11) an expression for midpoint wavenumber\index{wavenumber!midpoint}
\begin{eqnarray}\text{\textit{$k$}}_{\text{\textit{$x$}}}& = & \frac{\omega }{\text{\textit{$c$}}} \left(\sin  \phi _{\text{\textit{$i$}}}-\sin  \phi
_{\text{\textit{$r$}}}\right)\nonumber\\
& = & \frac{2\omega }{\text{\textit{$c$}}} \sin  \alpha  \cos  \gamma ,\end{eqnarray}
and from (3.12), for offset wavenumber,\index{wavenumber!offset}
\begin{eqnarray}\text{\textit{$k$}}_{\text{\textit{$h$}}} & = & -\frac{\omega }{\text{\textit{$c$}}} \left(\sin  \phi _{\text{\textit{$i$}}}+\sin  \phi
_{\text{\textit{$r$}}}\right)\nonumber\\
& = & -\frac{2\omega }{\text{\textit{$c$}}} \sin  \gamma  \cos  \alpha .\end{eqnarray}

From these formulas it follows that dip angle\index{angle!dip} $\alpha $ depends on the ratio \(k_x/k_z\),
\begin{equation}\tan  \alpha  = -\text{\textit{$ $}}\text{\textit{$k$}}_{\text{\textit{$x$}}}/\text{\textit{$k$}}_{\text{\textit{$z$}}},\end{equation}
while angle of incidence\index{angle!incident} $\gamma $ depends on \(k_h/k_z\):
\begin{equation}\tan  \gamma  =\text{\textit{$ $}}\text{\textit{$k$}}_{\text{\textit{$h$}}}/\text{\textit{$k$}}_{\text{\textit{$z$}}},\end{equation}
or
\begin{equation}\cos  \gamma  =\frac{1}{\sqrt{1+\text{\textit{$k$}}_{\text{\textit{$h$}}}{}^2/\text{\textit{$k$}}_{\text{\textit{$z$}}}{}^2}}.\end{equation}

Equation (3.32) is the relation needed to evaluate \(\left.dk_h\right/d\gamma \). We have
\begin{equation}\frac{d\text{\textit{$k$}}_{\text{\textit{$h$}}}}{d \gamma } = \text{\textit{$ $}}\frac{\text{\textit{$k$}}_{\text{\textit{$z$}}}}{\cos
 ^2\gamma }.\end{equation}

Hence, (3.22) becomes
\begin{equation}\text{\textit{$\mathit{r}$}}\left(\text{\textit{$k_x$}},\text{\textit{$k_z$}}\text{\textit{$,$}}\gamma \right) =\frac{-\text{\textit{$c$}}^2}{4\pi
} \text{\textit{$D$}}\left(\text{\textit{$k$}}_{\text{\textit{gx}}}|\text{\textit{$k$}}_{\text{\textit{sx}}};\omega \right) \frac{1}{|\omega |\sqrt{-i
\omega }}\frac{\text{\textit{$k$}}_{\text{\textit{gz}}}\text{\textit{$k$}}_{\text{\textit{sz}}}}{ \cos  ^2\gamma } .\end{equation}

Equations (3.17), (3.18), and (3.19) for $\omega $, \(k_{\text{\textit{sz}}}\), and \(k_{\text{\textit{gz}}}\) can be expressed in terms of
\(k_x\), \(k_z\), and $\gamma $ as
\begin{align}\omega &=-\text{\textit{$k_z$}}\frac{\text{\textit{$c$}}}{2}\sqrt{1+\text{\textit{$k$}}_{\text{\textit{$x$}}}{}^2/\text{\textit{$k$}}_{\text{\textit{$z$}}}{}^2}
\sec  \gamma ,\\
\text{\textit{$k$}}_{\text{\textit{sz}}}&=-\frac{\text{\textit{$k$}}_{\text{\textit{$z$}}}}{2} \left(1-\frac{\text{\textit{$k$}}_{\text{\textit{$x$}}}\tan
 \gamma }{\text{\textit{$k$}}_{\text{\textit{$z$}}}}\right),\end{align}
and
\begin{equation}\text{\textit{$k$}}_{\text{\textit{gz}}}=\frac{\text{\textit{$k$}}_{\text{\textit{$z$}}}}{2} \left(1+\frac{\text{\textit{$k$}}_{\text{\textit{$x$}}}\tan
 \gamma }{\text{\textit{$k$}}_{\text{\textit{$z$}}}}\right) .\end{equation}

Equation (3.35) gives a one-to-one mapping between \textit{f--k} data points and the \hbox{\textit{k}-space} migrated image at a particular angle of incidence.
The amplitude factor in (3.35) should not be taken seriously. The most one can say is that, with the above amplitude, the integral over incidence
angle of the angle-dependent image function \(\text{\textit{$\mathit{r}$}}\left(\text{\textit{$k$}}_x,k_z,\gamma \right)\) is the image function\index{image function}
\(R\left(k_x,k_z\right)\), as defined in equation (3.16).

\subsection{Amplitude-preserving prestack f--k migration in 2.5 dimensions}

Guided by the form (3.35), let's derive an amplitude-faithful expression for 2-D prestack \textit{f--k} migration\index{migration!prestack amplitude preserving 2.5D $f{-}k$}. We seek an expression
of the form
\begin{equation}\text{\textit{$D$}}\left(\text{\textit{$k$}}_{\text{\textit{$x$}}},\text{\textit{$k$}}_{\text{\textit{$h$}}},\omega \right)\text{
 }= \text{\textit{$R$}}\left(\text{\textit{$k$}}_{\text{\textit{$x$}}},\text{\textit{$k_z$}}\text{\textit{$,$}}\gamma \right) \mathcal{A}(\alpha
,\gamma ,\omega ),\end{equation}
with $\alpha $ the dip angle\index{angle!dip}, defined as
\begin{equation}\tan  \alpha  =- \text{\textit{$k_x$}}/\text{\textit{$k_z$}},\end{equation}
$\gamma $ the incident angle\index{angle!incident}, given by
\begin{equation}\tan  \gamma  = \text{\textit{$k_h$}}/\text{\textit{$k_z$}},\end{equation}
and frequency $\omega $ given by
\begin{equation}\omega =-\frac{\text{\textit{$k_zc$}}}{2}\sqrt{\left(1+\text{\textit{$k_x$}}{}^2/\text{\textit{$k_z$}}{}^2\right)\left(1+\text{\textit{$k_h$}}{}^2/\text{\textit{$k_z$}}{}^2\right)},\end{equation}
also expressible as (3.28).

To figure out what $\mathcal{A}$ must be, look at a single dipping reflector. Invoking the principle of superposition,\index{superposition!principle of} we claim that a formula
that can migrate a dipping reflector can migrate anything. Spatially, given a reflector at \(z=z_0+x \tan  \alpha _0\), the formula for \textit{R} is
\begin{equation}\text{\textit{$R$}}(\text{\textit{$x$}},\text{\textit{$z$}}\text{\textit{$,$}}\gamma ) =\text{\textit{$R$}}_0(\gamma )\delta \left(\left(\text{\textit{$z$}}-\text{\textit{$z$}}_0\right)\cos
 \alpha _0-\text{\textit{$x$}}\text{\textit{$ $}}\sin  \alpha _0\right).\end{equation}

In terms of wavenumber, \textit{R} is
\begin{eqnarray}
\text{\textit{$R$}}\left(\text{\textit{$k$}}_{\text{\textit{$x$}}},\text{\textit{$k$}}_{\text{\textit{$z$}}}\text{\textit{$,$}}\gamma
\right) & = & \text{\textit{$R$}}_0(\gamma )\int \int d\text{\textit{$x$}}d\text{\textit{$z$}}\text{\textit{$ $}}\delta \left(\left(\text{\textit{$z$}}-\text{\textit{$z$}}_0\right)\cos
 \alpha _0-\text{\textit{$x$}}\text{\textit{$ $}}\sin  \alpha _0\right)e^{-i\left(\text{\textit{$k_xx$}}+\text{\textit{$k_zz$}}\right)} \nonumber\\
& = & 2\pi \text{\textit{$R$}}_0(\gamma )e^{-i \text{\textit{$k_z$}}\text{\textit{$z$}}_0}\delta \left(\text{\textit{$k_x$}}\cos  \alpha _0+\text{\textit{$k_z$}}\sin
 \alpha _0\right).
\end{eqnarray}

The delta function confirms what we should already know for a dipping layer, namely (see equation (3.31)) tan \(\alpha _0\) = \(-k_x/k_z\).  By
definition, tangent of incident angle is given by (3.32), or tan $\gamma $ = \(k_h/k_z\). Since dip is now fixed at \(\alpha _0\), the equation
(3.42) for frequency $\omega $ reduces to
\begin{equation}\omega =-\frac{\text{\textit{$k_zc$}}}{2}\sec  \alpha _0\sqrt{1+\text{\textit{$k_h$}}{}^2/\text{\textit{$k_z$}}{}^2},\end{equation}
or, inversely,
\begin{equation}\text{\textit{$k_z$}}=-\text{Sign}(\omega )\sqrt{\frac{4\omega ^2}{\text{\textit{$c$}}^2}\cos ^2\alpha _0-\text{\textit{$k_h$}}{}^2}.\end{equation}

The space--time formula for the reflection response \textit{D} can be determined from the method of images. In two dimensions, the impulse response
is a Hankel function\index{Hankel function} (see, e.g., \cite{tufte_1997}), corresponding to line sources and receivers. Typically, however, we want point sources
and receivers, even if the experiment is otherwise two-dimensional. In three dimensions, the reflected impulse response is a delta function
\begin{equation}\text{\textit{$D$}}\left(\text{\textit{$x_m$}},\text{\textit{$x_h$}},\text{\textit{$t$}}\right)\text{  }= \frac{\text{\textit{$R$}}_0(\gamma
) \delta (\text{\textit{$t$}}-\text{\textit{$r$}}/\text{\textit{$c$}})}{\text{\textit{$r$}}},\end{equation}
with \textit{r} the total raypath\index{raypath} length (see exercise 3.1),
\begin{equation}\text{\textit{$r$}}= 2\cos  \alpha _0\sqrt{\text{\textit{$x_h$}}{}^2+\text{\textit{$z_m$}}{}^2},\end{equation}
\(z_m\) the reflector depth directly below the midpoint \(x_m\),
\begin{equation}\text{\textit{$z_m$}}\equiv \text{\textit{$z$}}_0+\text{\textit{$x_m$}}\text{\textit{$ $}}\tan  \alpha _0,\end{equation}
and $\gamma $ the ray-theoretical angle of incidence (see exercise 3.1)
\begin{equation} \gamma =-\text{arctan}\left(\frac{\text{\textit{$x_h$}}\text{  }}{\text{\textit{$z_m$}}}\right).\end{equation}

The amplitude decay in the impulse response of 1/\textit{r} is referred to as spherical divergence.

A minor problem in linking \textit{R} to \(D\) is that the 2-D wave equation anticipates cylindrical, not spherical divergence. We
address this problem by multiplying the 3-D data by the square root of time:
\begin{eqnarray}\text{\textit{$D$}}^{(\text{{tdc}})}\left(\text{\textit{$x_m$}},\text{\textit{$x_h$}},
\text{\textit{$t$}}\right) & \equiv &
\sqrt{\text{\textit{$t$}}}\text{\textit{$D$}}\left(\text{\textit{$x_m$}},\text{\textit{$x_h$}},\text{\textit{$t$}}\right)\nonumber\\
& = &  \frac{\text{\textit{$R$}}_0(\gamma
) \delta (\text{\textit{$t$}}-\text{\textit{$r$}}/\text{\textit{$c$}})}{\sqrt{\text{\textit{$r c$}}}}.\end{eqnarray}

We expect to be able, at least asymptotically, to relate \textit{R} and \(D^{(\text{{tdc}})}\) (tdc standing for time divergence corrected) with a 2-D mapping of the form (3.39). Although
there is also a phase difference between the 2-D Hankel function and the 3-D impulse response, we hope to include that in the coefficient $\mathcal{A}$.

To confirm the form (3.39), we must Fourier transform \(D^{(\text{{tdc}})}\) over \(x_m\), \(x_h\), and \textit{t, }and compare to \(R\left(k_x,k_z,\gamma
\right)\).
\begin{equation}\text{\textit{$D$}}^{(\text{{tdc}})}\left(\text{\textit{$k_m$}},\text{\textit{$k_h$}},\omega \right)\text{  }= \int d\text{\textit{$x_m$}}\int
d\text{\textit{$x_h$}}\frac{\text{\textit{$R$}}_0(\gamma ) e^{i \left(\omega  \text{\textit{$r/c$}}\text{\textit{$-$}}\text{\textit{$k_hx_h$}}\text{\textit{$-$}}\text{\textit{$\text{\textit{$k$}}_mx_m$}}\right)}}{\text{\textit{$\sqrt{\text{\textit{$r$}}
c}$}}}.\end{equation}

Whereas \textit{R} is a function of \(k_x\), \(k_z\), and $\gamma $, \(D^{(\text{{tdc}})}\) is a function of \(k_m, k_h\), and $\omega $.

Perform the integral over offset using the stationary phase approximation\index{stationary phase approximation}. We seek the stationary point of the phase argument
\begin{eqnarray}\text{\textit{$f$}}\left(\text{\textit{$x_h$}}\right) & = & \omega  \text{\textit{$r / c$}} -\text{\textit{$ $}}\text{\textit{$k_hx_h$}}\nonumber\\
& = & \frac{2\omega  }{\text{\textit{$c$}}}\cos  \alpha _0 \sqrt{\text{\textit{$x_h$}}{}^2+\text{\textit{$z_m$}}{}^2}-\text{\textit{$ $}}\text{\textit{$k_hx_h$}}\text{\textit{$,$}}\end{eqnarray}
i.e. the point at which its derivative,
\begin{equation}\text{\textit{$f'$}}\left(\text{\textit{$x_h$}}\right)=\frac{2\omega  \cos  \alpha _0\text{\textit{$ $}}\text{\textit{$x_h$}}}{\text{\textit{$
$}}\text{\textit{$c$}}\text{\textit{$ $}}\sqrt{\text{\textit{$x_h$}}{}^2+\text{\textit{$z_m$}}{}^2}} -\text{\textit{$ $}}\text{\textit{$k_h$}}\text{\textit{$,$}}\end{equation}
is zero. The stationary point is
\begin{equation}\text{\textit{$ $}}\text{\textit{$\underset{{}^{\wedge}}{x}{}_h$}} =\text{\textit{$z_m$}}\frac{\text{\textit{$ $}}\text{\textit{$k_h$}}\text{\textit{
 }}\text{\textit{$c$}}\ \text{\textit{$ $}}\text{Sign}\text{\textit{$($}}\omega )}{\text{  }\sqrt{4\omega ^2\cos ^2\alpha _0/\text{\textit{$c$}}^2-\text{\textit{$k_h$}}{}^2}}\text{\textit{$
$}}=\text{\textit{$-z_m$}}\frac{\text{\textit{$ k_h $}}}{\text{\textit{$k_z$}}},\end{equation}
at which
\begin{equation}\sqrt{\text{\textit{$\underset{{}^{\wedge}}{x}{}_h$}}{}^2+\text{\textit{$z_m$}}{}^2} =\text{\textit{$\frac{2|\omega |\cos  \alpha
_0 \text{\textit{$z_m$}}}{\text{\textit{$c$}} \sqrt{4\omega ^2\cos ^2\alpha _0/\text{\textit{$c$}}^2-\text{\textit{$k_h$}}{}^2}}$}}\text{\textit{$
$}}=-\frac{2 \omega \text{  }\cos  \alpha _0 \text{\textit{$z_m$}}}{\text{\textit{$c k_z$}}}\text{\textit{$,$}}\end{equation}
\begin{equation}\text{\textit{$\underset{{}^{\wedge}}{\text{\textit{$r$}}}$}}=-4\cos ^2 \alpha _0\text{\textit{$\frac{\omega  \text{\textit{$z_m$}}}{\text{\textit{$c
k_z$}}}$}},\end{equation}
and
\begin{equation}\text{\textit{$f$}}\left(\text{\textit{$\underset{{}^{\wedge}}{x}{}_h$}}\right)=\text{\textit{$ $}}\text{\textit{$-$}}\text{\textit{$k$}}_{\text{\textit{$z$}}}\text{\textit{$z_m$}}\text{\textit{$,$}}\end{equation}
the expressions in terms of vertical wavenumber \(k_z\) following from (3.46).

Note also that the ray-theoretical expression (3.50) for angle of incidence becomes the wave-theoretical expression (3.41):
\begin{equation}\underset{{}^{\wedge}}{\gamma }=\underset{{}^{\wedge}}{\gamma }\left(\text{\textit{$k_m$}},\text{\textit{$k_h$}},\omega \right) -\text{arctan}\left(\text{\textit{$
$}}\text{\textit{$\underset{{}^{\wedge}}{x}{}_h/z_m$}}\right)=\text{arctan}\left(\text{\textit{$ $}}\text{\textit{$k_h /k_z$}}\right).\end{equation}

This allows us to write the relation between $\omega $ and \(k_z\) in terms of \(\underset{{}^{\wedge}}{\gamma }\):
\begin{equation}\text{\textit{$k_z$}}=-\frac{2\omega }{\text{\textit{$c$}}}\cos  \alpha _0 \cos  \underset{{}^{\wedge}}{\gamma }.\end{equation}

The second derivative of the phase argument is
\begin{equation}\text{\textit{$f^{\prime\prime}$}}\left(\text{\textit{$x_h$}}\right)=\frac{2\omega  \cos  \alpha _0\text{\textit{$ $}}\text{\textit{$z_m$}}{}^2
}{\text{\textit{$ $}}\text{\textit{$c$}}\text{\textit{$ $}}\left.\text{\textit{$($}}\text{\textit{$x_h$}}{}^2+\text{\textit{$z_m$}}{}^2\right){}^{3/2}},\end{equation}
with stationary value
\begin{equation}\text{\textit{$f^{\prime\prime}$}}\left(\text{\textit{$\underset{{}^{\wedge}}{x}{}_h$}}\right)=\frac{2\omega  \cos  \alpha _0\text{\textit{$
$}}\cos ^3 \underset{{}^{\wedge}}{\gamma }}{\text{\textit{$ c z_m$}}}=-\frac{\text{\textit{$k_z$}}\cos ^2 \underset{{}^{\wedge}}{\gamma }}{\text{\textit{$z_m$}}}.\end{equation}

The stationary phase approximation\index{stationary phase approximation} to the offset integral thus yields
\begin{eqnarray}
\text{\textit{$D$}}^{(\text{{tdc}})}\left(\text{\textit{$k_m$}},\text{\textit{$k_h$}},\omega \right)\text{  } & = & \int d\text{\textit{$x_m$}}e^{-i
\text{\textit{$\text{\textit{$k$}}_m$}}\text{\textit{$x_m$}}} \sqrt{\frac{2\pi  i}{\text{\textit{$\underset{{}^{\wedge}}{\text{\textit{$r$}}}$}}
\text{\textit{$cf^{\prime\prime}$}}\left(\text{\textit{$\underset{{}^{\wedge}}{x}{}_h$}}\right)}}\text{\textit{$R$}}_0\left(\underset{{}^{\wedge}}{\gamma
}\right) e^{i \text{\textit{$f$}}\left(\text{\textit{$\underset{{}^{\wedge}}{x}{}_h$}}\right)} \nonumber\\
& = &  \int d\text{\textit{$x_m$}}\sqrt{\frac{\pi  i }{2\omega  }}\frac{\text{\textit{$R$}}_0\left(\underset{{}^{\wedge}}{\gamma }\right)}{\cos  \underset{{}^{\wedge}}{\gamma
} \cos \text{\textit{$ $}}\alpha _0} e^{-i \left(\text{\textit{$k_zz_m$}}+\text{\textit{$\text{\textit{$k$}}_mx_m$}}\text{\textit{$)$}}\right.}.
\end{eqnarray}

Comparing with equation (3.44), we invoke the principle of superposition\index{superposition!principle of} to assert that for \textit{R} a general function of \(k_m\), \(k_z\), and $\gamma$,
\begin{equation}\text{\textit{$D$}}^{(\text{{tdc}})}\left(\text{\textit{$k$}}_{\text{\textit{$x$}}},\text{\textit{$k$}}_{\text{\textit{$h$}}},\omega
\right)\text{  }=\sqrt{\frac{\pi  i \text{\textit{$ $}}}{2\omega  }}\frac{\text{\textit{$R$}}\left(\text{\textit{$k$}}_{\text{\textit{$m$}}},\text{\textit{$k_z$}}\text{\textit{$,$}}\, \gamma
\right)}{\cos  \gamma }.\end{equation}

This is of the form (3.39), with
\begin{equation} \mathcal{A}^{(\text{{tdc}})}(\alpha ,\gamma ,\omega ) \equiv \frac{1}{\cos  \gamma }\sqrt{\frac{\pi  i}{2\omega  }}.\end{equation}

This expression should (at least asymptotically) preserve reflection amplitude and phase. It does have the same phase as the more naive expression
(3.35), though of course a different amplitude. An equivalent 2.5-D expression is derived in \nobreak Chapter 9 via scattering theory.\index{scattering theory} It differs from
(3.65) in that there, the 2.5-D divergence correction (3.51) is performed after the migration. The difference is superficial: the divergence
correction may be applied either before migration as \(\sqrt{\text{\textit{$t$}}}\), or after migration as \(\sqrt{\text{\textit{$z$}}}\) with
a slightly altered amplitude coefficient.

\section{Prestack slant-stack migration in 2.5 dimensions}

Expressing prestack migration in terms of slant stacks or Radon
transforms\index{Radon transform} is surprisingly easy. Start with
the amplitude-preserving 2.5-D formula (3.65), with the
\textit{k}-space representation of \textit{R} replaced by
\begin{equation}\text{\textit{$R$}}^{(2.5\text{D})}\left(\text{\textit{$q$}}_{\text{\textit{$x$}}},\text{\textit{$q$}}_{\text{\textit{$h$}}},\text{\textit{$k_z$}}\right)\text{
 }= \text{\textit{$R$}}^{(2.5\text{D})}\left(\text{\textit{$k$}}_{\text{\textit{$x$}}},\text{\textit{$k_z$}}\text{\textit{$,$}}\, \gamma \right),\end{equation}
with \(q_x=k_x/k_z=-\tan  \alpha \) and \(q_h=k_h/k_z=\tan  \gamma \). In this representation, \textit{R} is the Fourier--Radon transform
\begin{equation}\text{\textit{$R$}}^{(2.5\text{D})}\left(\text{\textit{$q$}}_{\text{\textit{$x$}}},\text{\textit{$q$}}_{\text{\textit{$h$}}},\text{\textit{$k_z$}}\right)\text{
 }= \int \int d\text{\textit{$x_m$}}d\text{\textit{$z$}}\text{\textit{$ $}}\text{\textit{$R$}}^{(2.5\text{D})}\left(\text{\textit{$x_m$}},\text{\textit{$z$}}-\text{\textit{$q$}}_{\text{\textit{$x$}}}\text{\textit{$x_m$}}\text{\textit{$,$}}\, \gamma
\right)e^{-i \text{\textit{$k_z$}}\text{\textit{$z$}}} .\end{equation}

Equation (3.65) becomes
\begin{equation}\text{\textit{$R$}}^{(2.5\text{D})}\left(\text{\textit{$q$}}_{\text{\textit{$x$}}},\text{\textit{$q$}}_{\text{\textit{$h$}}},\text{\textit{$k_z$}}\right)
=\text{\textit{$D$}}^{(2.5\text{D})}\left(\text{\textit{$k$}}_{\text{\textit{$x$}}},\text{\textit{$k$}}_{\text{\textit{$h$}}},\omega \right) \sqrt{\frac{-2
i \omega }{\pi \left(1+\text{\textit{$q_h$}}{}^2\right)}} .\end{equation}

In terms of the new parameters, the relation between \(k_z\) and $\omega $ is
\begin{equation}\omega =-\text{\textit{$\frac{\text{\textit{$k_z$}}c}{2}$}}\sqrt{\left(1+\text{\textit{$q_x$}}{}^2\right)\left(1+\text{\textit{$q_h$}}{}^2\right)}.\end{equation}

This operation is a 2-D generalized slant stack, in which the time coordinate is not only shifted but compressed.

\section{Prestack ray-theoretical migration}

Given almost any source--receiver configuration, it is possible to perform a migration using ray-theoretical concepts. If rays can be computed
from source and receiver locations to arbitrary points in the subsurface, then for a given source and receiver, a map can be generated of traveltimes
from source to all possible reflection points to receiver. An event at traveltime $\tau $ can be broadcast to
all possible reflection points consistent with traveltime $\tau$ using the method of \nobreak Section~2.2.4. \nobreak Combining the data from many sets of sources and receivers, we can expect an
event to build constructively along actual reflector locations, with destructive interference attenuating the response at locations where no reflector
exists. The principal requirement, other than the availability of rays, is that the source--receiver configuration result in dense coverage of
the image space. Typically, ray-theoretical migration uses a common-offset data set, though other configurations are certainly possible.

For constant velocity, if amplitude is not an issue we can use the amplitude-challenged 2-D \textit{f--k} formula (3.10) to derive an asymptotically
equivalent Kirchhoff migration algorithm\index{migration!Kirchhoff}. We return the data to space--time, obtaining (after some rearranging)


Since \(k_{\text{\textit{gz}}}\) is a function of \(k_{\text{\textit{gx}}}\) and $\omega $, and \(k_{\text{\textit{sz}}}\) is a function of \(k_{\text{\textit{sx}}}\)
and $\omega $ (see equations (3.5) and (3.7)), the phases in the \(k_{\text{\textit{gx}}}\) and \(k_{\text{\textit{sx}}}\) integrals are nonlinear
in their arguments. Exact solutions to these integrals are possible, but it is more convenient to seek an asymptotic solution through the stationary
phase approximation. Invoking Appendix F, the \(k_{\text{\textit{gx}}}\) and \(k_{\text{\textit{sx}}}\) integrals can be approximately evaluated
to~be
\begin{equation}\int d\text{\textit{$k$}}_{\text{\textit{sx}}}\text{  }e^{-i\left(\text{\textit{$k$}}_{\text{\textit{sz}}}\text{\textit{$z$}}+\text{\textit{$k$}}_{\text{\textit{sx}}}\left(\text{\textit{$x$}}\text{\textit{$-$}}\text{\textit{$x$}}_{\text{\textit{$s$}}}\right)\text{\textit{$)$}}\right.}\simeq
e^{-i \omega  \left.\text{\textit{$r$}}_{\text{\textit{$s$}}}\right/\text{\textit{$c$}}}\sqrt{\frac{2\pi  i \omega  \text{\textit{$z$}}^2}{\text{\textit{$c$}}\text{\textit{$
$}}\text{\textit{$r$}}_{\text{\textit{$s$}}}{}^3}},\end{equation}
and
\begin{equation}\int d\text{\textit{$k$}}_{\text{\textit{gx}}}\text{  }e^{i\left(\text{\textit{$k$}}_{\text{\textit{gz}}}\text{\textit{$z$}}+\text{\textit{$k$}}_{\text{\textit{gx}}}\left(\text{\textit{$x$}}\text{\textit{$-$}}\text{\textit{$x$}}_{\text{\textit{$g$}}}\right)\text{\textit{$)$}}\right.}\simeq
e^{-i \omega  \left.\text{\textit{$r$}}_{\text{\textit{$g$}}}\right/\text{\textit{$c$}}}\sqrt{\frac{2\pi  i \omega  \text{\textit{$z$}}^2}{\text{\textit{$c$}}\text{\textit{$
$}}\text{\textit{$r$}}_{\text{\textit{$g$}}}{}^3}},\end{equation}
with \(r_s\) the distance from source \(\left(x_s,0\right)\) to image point (\textit{x},\, \textit{z})
\begin{equation}\text{\textit{$r$}}_{\text{\textit{$s$}}}=\sqrt{\text{\textit{$z$}}^2+\left(\text{\textit{$x$}}-\text{\textit{$x$}}_{\text{\textit{$s$}}}\right){}^2}\end{equation}
and \(r_g\) the distance from image point (\textit{x}, \textit{z}) to receiver \(\left(x_g,0\right)\)
\begin{equation}\text{\textit{$r$}}_{\text{\textit{$g$}}}=\sqrt{\text{\textit{$z$}}^2+\left(\text{\textit{$x$}}-\text{\textit{$x$}}_{\text{\textit{$g$}}}\right){}^2},\end{equation}
as depicted in Figure 3.3. Thus, setting
\begin{equation}\text{\textit{$r$}}=\text{\textit{$r$}}_{\text{\textit{$g$}}}+\text{\textit{$r$}}_{\text{\textit{$s$}}},\end{equation}
the image function \textit{R} becomes
\begin{equation}\text{\textit{$R$}}(\text{\textit{$x,z$}}) = \frac{\text{\textit{$z$}}^2}{(2\pi )^2\text{\textit{$c$}}} \int d\text{\textit{$x$}}_{\text{\textit{$g$}}}\int
d\text{\textit{$x$}}_{\text{\textit{$s$}}}\int d\text{\textit{$t$}}\text{\textit{$ $}}\frac{\text{\textit{$D$}}\left(\text{\textit{$x$}}_{\text{\textit{$g$}}}|\text{\textit{$x$}}_{\text{\textit{$s$}}};\text{\textit{$t$}}\right)}{\left(\text{\textit{$r$}}_{\text{\textit{$s$}}}\text{\textit{$r$}}_{\text{\textit{$g$}}}\right){}^{3/2}}\int
d\omega  \sqrt{-i \omega } e^{i \omega (\text{  }\text{\textit{$t$}}\text{\textit{$-$}}\text{\textit{$r$}}/\text{\textit{$c$}})}. \end{equation}

The frequency integral can be evaluated as the half-derivative filter\index{filter!half-derivative}
\begin{equation}\int d\omega  \sqrt{-i \omega } e^{i \omega (\text{  }\text{\textit{$t$}}\text{\textit{$-$}}\text{\textit{$r$}}/\text{\textit{$c$}})}=
2\pi  \text{\textit{$d_{h+}$}}( \text{\textit{$t$}}\text{\textit{$-$}}\text{\textit{$r$}}/\text{\textit{$c$}}).\end{equation}

The time integral is a convolution of \(d_{h+}\) with the data function:
\begin{equation}\text{\textit{$R$}}(\text{\textit{$x,z$}})= \frac{\text{\textit{$z$}}^2}{(2\pi )\text{\textit{$c$}}} \int d\text{\textit{$x$}}_{\text{\textit{$g$}}}\int
d\text{\textit{$x$}}_{\text{\textit{$s$}}}\frac{\left[\text{\textit{$d_{h+}*D$}}\text{\textit{$]$}}\right.\left(\text{\textit{$x$}}_{\text{\textit{$g$}}}|\text{\textit{$x$}}_{\text{\textit{$s$}}};\text{\textit{$t$}}=\text{\textit{$r$}}/\text{\textit{$c$}}\right)}{\left(\text{\textit{$r$}}_{\text{\textit{$s$}}}\text{\textit{$r$}}_{\text{\textit{$g$}}}\right){}^{3/2}}.\end{equation}

As was the case for the \textit{f--k} formula from which this expression is derived, one should not take either the amplitude or the phase too seriously.


According to this equation, full Kirchhoff migration\index{migration!Kirchhoff} is a weighted summation of the data over all source and receiver positions, along a trajectory
of traveltimes corresponding to raypaths from the source position to the image point to the receiver point. This trajectory function is plotted
as a function of \(x_g\) and \(x_s\) in Figure 3.4.

The formula (3.84) does not imply that the summation \textit{must} be over all source and receiver positions. One can pluck out a subset of the
complete data set, corresponding to a line on the surface depicted in Figure 3.4. An example would be a constant-offset section. With \(x_h=\left.\left(x_g-x_s\right)\right/2\),
and \(x_m=\left.\left(x_g+x_s\right)\right/2\), we can write (3.84) as
\begin{equation}
\text{\textit{$R$}}\left(\tilde{\boldsymbol{x}}\right) = \int d\text{\textit{$x$}}_{\text{\textit{$h$}}}\text{\textit{$\mathcal{R}$}}\left(\tilde{\boldsymbol{x}},\text{\textit{$x_h$}}\right),\end{equation}
with
\begin{equation}\text{\textit{$\mathcal{R}$}}\left(\tilde{\boldsymbol{x}},\,\text{\textit{$x$}}_{\text{\textit{$h$}}}\right)
= \frac{\text{\textit{$z$}}^2}{(2\pi )\text{\textit{$c$}}} \int d\text{\textit{$x$}}_{\text{\textit{$m$}}}\frac{\text{\textit{$d_{h+}$}}\text{\textit{$*$}}\text{\textit{$D$}}\left(\text{\textit{$x$}}_{\text{\textit{$g$}}}|\text{\textit{$x$}}_{\text{\textit{$s$}}};\text{\textit{$t$}}=\text{\textit{$r$}}/\text{\textit{$c$}}\right)}{\left(\text{\textit{$r$}}_{\text{\textit{$s$}}}\text{\textit{$r$}}_{\text{\textit{$g$}}}\right){}^{3/2}}.\end{equation}

The trajectory function for this subset is
\begin{equation}\text{\textit{$c t$}}=\text{\textit{$r$}}=\sqrt{\text{\textit{$z$}}^2+\left(\Delta \text{\textit{$x$}}-\text{\textit{$x_h$}}\right){}^2}+\sqrt{\text{\textit{$z$}}^2+\left(\Delta
\text{\textit{$x$}}+\text{\textit{$x_h$}}\right){}^2},\end{equation}
with \(x_m-x\equiv \Delta \text{\textit{$x$}}\). This expression can be solved for \textit{z}, obtaining
\begin{equation}\text{\textit{$z$}}=\frac{\text{\textit{$c t$}}}{2\text{\textit{$ $}}}\sqrt{\left(1-4 \frac{\Delta \text{\textit{$x$}}^2}{\text{\textit{$c$}}^2\text{\textit{$t$}}^2}\right)\left(1-4
\frac{\text{\textit{$x_h$}}{}^2}{\text{\textit{$c$}}^2\text{\textit{$t$}}^2}\right)}.\end{equation}

Substituting this expression for \textit{z}, the ray lengths \(r_s\) and \(r_g\) can be rewritten as
\begin{equation}\text{\textit{$r$}}_{\text{\textit{$s$}}}=\frac{\text{\textit{$c t$}}}{2}-2\frac{\Delta \text{\textit{$x$}}\text{\textit{$ $}}\text{\textit{$x_h$}}}{\text{\textit{$c
t$}}},\end{equation}
\begin{equation}\text{\textit{$r$}}_{\text{\textit{$g$}}}=\frac{\text{\textit{$c t$}}}{2}+2\frac{\Delta \text{\textit{$x$}}\text{\textit{$ $}}\text{\textit{$x_h$}}}{\text{\textit{$c
t$}}}.\end{equation}

The angles made by the incident and reflected rays with the vertical are given by the formulas
\begin{equation}\sin  \phi _{\text{\textit{$i$}}}=\frac{\text{\textit{$x_h$}}-\Delta \text{\textit{$x$}}\text{\textit{$ $}}}{\text{\textit{$r$}}_{\text{\textit{$s$}}}},\text{
  }\cos  \phi _{\text{\textit{$i$}}}=\frac{\text{\textit{$z$}}}{\text{\textit{$r$}}_{\text{\textit{$s$}}}},\end{equation}
and
\begin{equation}\sin  \phi _{\text{\textit{$r$}}}=\frac{\text{\textit{$x_h$}}+\Delta \text{\textit{$x$}}\text{\textit{$ $}}}{\text{\textit{$r$}}_{\text{\textit{$g$}}}},\ \cos
 \phi _{\text{\textit{$r$}}}=\frac{\text{\textit{$z $}}}{\text{\textit{$r$}}_{\text{\textit{$g$}}}}.\end{equation}

We have defined these angles to be positive when \(x_g>x\) and \(x_s<x\). Angle of incidence $\gamma $ is defined to be \(\gamma \equiv \left.\left(\phi
_{\text{\textit{$i$}}}+\phi _{\text{\textit{$r$}}}\right)\right/2\), so that
\begin{align}\cos  2\gamma &=\cos  \left(\phi _{\text{\textit{$i$}}}+\phi _{\text{\textit{$r$}}}\right)=\cos  \phi _{\text{\textit{$i$}}}\cos
\phi _{\text{\textit{$r$}}}-\sin  \phi _{\text{\textit{$i$}}} \sin  \phi _{\text{\textit{$r$}}}\nonumber\\
&=\frac{\text{\textit{$z$}}^2\text{\textit{$ $}}+\Delta \text{\textit{$x$}}^2\text{\textit{$-$}}\text{\textit{$x_h$}}{}^2}{\text{\textit{$r$}}_{\text{\textit{$s$}}}\text{\textit{$r$}}_{\text{\textit{$g$}}}}=\frac{1-8
\frac{\text{\textit{$ $}}\text{\textit{$x_h$}}{}^2}{\text{\textit{$c$}}^2\text{\textit{$ $}}\text{\textit{$t$}}^2}+16\frac{\Delta \text{\textit{$x$}}^2\text{\textit{$
$}}\text{\textit{$x_h$}}{}^2}{\text{\textit{$c$}}^4\text{\textit{$ $}}\text{\textit{$t$}}^4}}{1-16\frac{\Delta \text{\textit{$x$}}^2\text{\textit{$
$}}\text{\textit{$x_h$}}{}^2}{\text{\textit{$c$}}^4\text{\textit{$ $}}\text{\textit{$t$}}^4}},
\end{align}
or
\begin{equation}\cos  \gamma =\sqrt{\frac{1 + \cos  2\gamma }{2}}=\sqrt{\frac{1 -4 \frac{\text{\textit{$ $}}\text{\textit{$x_h$}}{}^2}{\text{\textit{$c$}}^2\text{\textit{$
$}}\text{\textit{$t$}}^2}}{1-16\frac{\Delta \text{\textit{$x$}}^2\text{\textit{$ $}}\text{\textit{$x_h$}}{}^2}{\text{\textit{$c$}}^4\text{\textit{$
$}}\text{\textit{$t$}}^4}}},\end{equation}
\begin{equation}\sin  \gamma =\frac{2\text{\textit{$ $}}\text{\textit{$x_h$}}}{\text{\textit{$c t$}}}\sqrt{\frac{1-4\frac{\Delta \text{\textit{$x$}}^2\text{\textit{$
$}}}{\text{\textit{$c$}}^2\text{\textit{$ $}}\text{\textit{$t$}}^2} }{1-16\frac{\Delta \text{\textit{$x$}}^2\text{\textit{$ $}}\text{\textit{$x_h$}}{}^2}{\text{\textit{$c$}}^4\text{\textit{$
$}}\text{\textit{$t$}}^4}}},\end{equation}
and
\begin{equation}\tan  \gamma =\frac{2\text{\textit{$ $}}\text{\textit{$x_h$}}}{\text{\textit{$c t$}}}\sqrt{\frac{1-4\frac{\Delta \text{\textit{$x$}}^2\text{\textit{$
$}}}{\text{\textit{$c$}}^2\text{\textit{$ $}}\text{\textit{$t$}}^2} }{1 -4 \frac{\text{\textit{$ $}}\text{\textit{$x_h$}}{}^2}{\text{\textit{$c$}}^2\text{\textit{$
$}}\text{\textit{$t$}}^2}}}.\end{equation}

This can also be written as
\begin{equation}\tan  \gamma =\frac{\text{\textit{$ x_h$}}}{\text{\textit{$z$}}}\left(1-4\frac{\Delta \text{\textit{$x$}}^2\text{\textit{$ $}}}{\text{\textit{$c$}}^2\text{\textit{$
$}}\text{\textit{$t$}}^2}\right).\end{equation}

Dip angle is \(\alpha =\left.\left(\phi _r-\phi _i\right)\right/2\), or
\begin{eqnarray}\cos  2\alpha &=&\cos  \left(\phi _{\text{\textit{$r$}}}+\phi _{\text{\textit{$i$}}}\right)=\cos  \phi _{\text{\textit{$i$}}}\cos
\phi _{\text{\textit{$r$}}}+\sin  \phi _{\text{\textit{$i$}}} \sin  \phi _{\text{\textit{$r$}}}\nonumber\\
&=&\frac{\text{\textit{$z$}}^2\text{\textit{$ $}}-\Delta \text{\textit{$x$}}^2+\text{\textit{$x_h$}}{}^2}{\text{\textit{$r$}}_{\text{\textit{$s$}}}\text{\textit{$r$}}_{\text{\textit{$g$}}}}=\frac{1-8
\frac{\text{\textit{$ $}}\Delta \text{\textit{$x$}}^2}{\text{\textit{$c$}}^2\text{\textit{$ $}}\text{\textit{$t$}}^2}+16\frac{\Delta \text{\textit{$x$}}^2\text{\textit{$
$}}\text{\textit{$x_h$}}{}^2}{\text{\textit{$c$}}^4\text{\textit{$ $}}\text{\textit{$t$}}^4}}{1-16\frac{\Delta \text{\textit{$x$}}^2\text{\textit{$
$}}\text{\textit{$x_h$}}{}^2}{\text{\textit{$c$}}^4\text{\textit{$ $}}\text{\textit{$t$}}^4}},
\end{eqnarray}
\begin{equation}\cos  \alpha =\sqrt{\frac{1 -4 \frac{\text{\textit{$ $}}\Delta \text{\textit{$x$}}^2}{\text{\textit{$c$}}^2\text{\textit{$ $}}\text{\textit{$t$}}^2}}{1-16\frac{\Delta
\text{\textit{$x$}}^2\text{\textit{$ $}}\text{\textit{$x_h$}}{}^2}{\text{\textit{$c$}}^4\text{\textit{$ $}}\text{\textit{$t$}}^4}}},\end{equation}
\begin{equation}\sin  \alpha =\frac{2\text{\textit{$ $}}\Delta \text{\textit{$x$}}}{\text{\textit{$c t$}}}\sqrt{\frac{1-4\frac{\text{\textit{$x_h$}}{}^2\text{\textit{$
$}}}{\text{\textit{$c$}}^2\text{\textit{$ $}}\text{\textit{$t$}}^2} }{1-16\frac{\Delta \text{\textit{$x$}}^2\text{\textit{$ $}}\text{\textit{$x_h$}}{}^2}{\text{\textit{$c$}}^4\text{\textit{$
$}}\text{\textit{$t$}}^4}}},\end{equation}
and
\begin{equation}\tan  \alpha =\frac{2\text{\textit{$ $}}\Delta \text{\textit{$x$}}}{\text{\textit{$c t$}}}\sqrt{\frac{1-4\frac{\text{\textit{$x_h$}}{}^2\text{\textit{$
$}}}{\text{\textit{$c$}}^2\text{\textit{$ $}}\text{\textit{$t$}}^2} }{1 -4 \frac{\text{\textit{$ $}}\Delta \text{\textit{$x$}}^2}{\text{\textit{$c$}}^2\text{\textit{$
$}}\text{\textit{$t$}}^2}}},\end{equation}
or, in mixed notation,
\begin{equation}\tan  \alpha =\frac{\text{\textit{$ $}}\Delta \text{\textit{$x$}}}{\text{\textit{$z$}}}\left(1-4\frac{\text{\textit{$x_h$}}{}^2\text{\textit{$
$}}}{\text{\textit{$c$}}^2\text{\textit{$ $}}\text{\textit{$t$}}^2}\right) .\end{equation}

Other subsets (such as constant source or receiver location) could similarly be extracted.

Either the full migration formula or its subsets can be generalized to variable velocity. If one is not too concerned with amplitudes, all one
need do is replace the straight-ray distances with those for variable-velocity rays. Constant or variable, the Kirchhoff formula lies in asymptopia.

\section{Prestack \textit{f--k} migration in three dimensions}
\sectionmark{Prestack {\rm f--k} migration in three dimensions}

Migration in \textit{f--k} space of a 3-D data set is a straightforward extension of 2-D migration. Without regard to the question as to whether a full, 5-D seismic experiment is practical, we want a migration operator of the form
\begin{equation}
\text{\textit{$R$}}(\boldsymbol{k},\gamma )=\text{\textit{$D$}}\left(\tilde{\boldsymbol{k}}_{\text{\textit{$g$}}}|
\tilde{\boldsymbol{k}}_{\text{\textit{$s$}}};\omega \right)\text{  }\mathcal{A}\left(\tilde{\boldsymbol{k}}_{\text{\textit{$g$}}}|\tilde{\boldsymbol{k}}_{\text{\textit{$s$}}};\omega \right),\end{equation}
where \textit{D} is a full 3-D surface data set, \textit{R} is the desired image function, and $\mathcal{A}$ is an amplitude-phase term. Data
are a function of the receiver wavenumber 2-vector \(\tilde{\boldsymbol{k}}_g=\left(k_{\text{\textit{gx}}}\text{\textit{$,$}}\,\text{\textit{$k$}}_{\text{\textit{gy}}}\right)\),
the source wavenumber 2-vector \(\tilde{\boldsymbol{k}}_s=\left(k_{\text{\textit{sx}}},k_{\text{\textit{sy}}}\right)\), and frequency $\omega $. The
wavenumbers \(\tilde{\boldsymbol{k}}_g\) and \(\tilde{\boldsymbol{k}}_s\) are the Fourier conjugates of receiver position \(\tilde{\boldsymbol{x}}_g=\left(x_g,y_g\right)\),
and source position \(\tilde{\boldsymbol{x}}_s=\left(x_s,y_s\right)\). The image function depends on image wavenumber \(\boldsymbol{k}\)
(the Fourier conjugate of image position \(\boldsymbol{x}\)), and also possibly on incident angle $\gamma $. (Conceivably, reflectivity\index{reflectivity} might also
depend on source--receiver azimuthal angle $\psi $, an additional parameter not present in the 2-D case. It appears in three dimensions because sources and
receivers are separated in two dimensions, and may assume many different relative orientations. With azimuth included, both the data and image functions
are 5-D, though for the models of reflection and transmission considered here, \textit{R} is supposed to be independent of azimuth.)
The amplitude-phase term $\mathcal{A}$ is also 5-D, though it is assumed to vary slowly with the spatial wavenumbers. We could
determine $\mathcal{A}$ in the same manner as in Section 3.2.1 or 3.2.2, but for now we will leave it open.

\subsection{Vertical wavenumbers}

The data are presumed to obey the wave equation, so source and receiver vertical wavenumber can be inferred from the horizontal wavenumbers. The
wave from source to reflector travels in the positive \textit{z} direction, so the source vertical wavenumber \(k_{\text{\textit{sz}}}\) will have
the same sign as frequency $\omega $:
\begin{equation}\text{\textit{$k$}}_{\text{\textit{sz}}}=+\frac{\omega }{\text{\textit{$c$}}}\text{  }\sqrt{1-\frac{\left(\text{\textit{$k$}}_{\text{\textit{sx}}}{}^2+\text{\textit{$k$}}_{\text{\textit{sy}}}{}^2\right)\text{\textit{$c$}}^2}{\omega
^2}}.\end{equation}

Waves from reflector to receiver travel in the negative $z$ direction, so \(k_{\text{\textit{gz}}}\) must have sign opposite to $\omega $:
\begin{equation}\text{\textit{$k$}}_{\text{\textit{gz}}}=- \frac{\omega }{\text{\textit{$c$}}} \sqrt{1-\frac{\left(\text{\textit{$k$}}_{\text{\textit{gx}}}{}^2+\text{\textit{$k$}}_{\text{\textit{gy}}}{}^2\right)\text{\textit{$c$}}^2}{\omega
^2}}.\end{equation}

The vertical wavenumbers \(k_{\text{\textit{gz}}}\) and \(k_{\text{\textit{sz}}}\) can be considered the last components of the 3-vectors
\begin{equation}\boldsymbol{k}_s=\left(\text{\textit{$k_{\text{sx}}$}},\text{\textit{$k_{\text{sy}}$}}\text{\textit{$,$}}\text{\textit{$k_{\text{sz}}$}}\right)\end{equation}
and
\begin{equation}\boldsymbol{k}_g=\left(\text{\textit{$k_{\text{gx}}$}},\text{\textit{$k_{\text{gy}}$}}\text{\textit{$,$}}\text{\textit{$k_{\text{gz}}$}}\right).\end{equation}

These vectors both have magnitude
\begin{equation}\text{\textit{$k_g$}}=\text{\textit{$k_s$}}=\frac{\omega }{\text{\textit{$c$}}}.\end{equation}

\subsection{Image function wavenumber}

We next define the 3-vector \pmb{\textit{k}} to be
\begin{equation}{\boldsymbol{k}}=\left(\text{\textit{$k_x$}},\text{\textit{$k_y$}}\text{\textit{$,$}}\text{\textit{$k_z$}}\right)=\text{\textit{$\pmb{\text{\textit{$k$}}}_g-\pmb{\text{\textit{$k$}}}_s.$}}\end{equation}

The wavenumber vector \pmb{ \textit{k}} points upward and bisects the source and receiver wavenumber vectors \(\boldsymbol{k}_s\) and \(\boldsymbol{k}_g\). {
}If a specular reflector exists at a reflection point, \pmb{ \textit{k}} is normal to it (see Figure H.1). This occurs because the Fourier transforms
have organized the data into components of particular dips. The element \(D\left(\tilde{\boldsymbol{k}}_g|\tilde{\boldsymbol{k}}_s;\omega
\right)\) contains all the data pertaining to reflectors normal to \pmb{ \textit{k}} = \(\boldsymbol{k}_g-\boldsymbol{k}_s\). A Fourier transform of the image
function \textit{R}(\pmb{ \textit{x}},\, $\gamma $) to \textit{R}(\pmb{ \textit{k}},\, $\gamma $) performs the same reflector dip decomposition on
the image, so we are justified in identifying the \pmb{ \textit{k}} of equation (3.109) with the image function wavenumber in (3.103).

The vertical component of image function wavenumber honors a 3-D double-square-root equation:
\begin{equation}\text{\textit{$k$}}_{\text{\textit{$z$}}}\equiv \text{\textit{$k$}}_{\text{\textit{gz}}}-\text{\textit{$k$}}_{\text{\textit{sz}}}\text{
 }=- \frac{\omega }{\text{\textit{$c$}}} \left( \sqrt{1-\frac{\left(\text{\textit{$k$}}_{\text{\textit{sx}}}{}^2+\text{\textit{$k$}}_{\text{\textit{sy}}}{}^2\right)\text{\textit{$c$}}^2}{\omega
^2}} +\sqrt{1-\frac{\left(\text{\textit{$k$}}_{\text{\textit{gx}}}{}^2+\text{\textit{$k$}}_{\text{\textit{gy}}}{}^2\right)\text{\textit{$c$}}^2}{\omega
^2}}\right).\end{equation}

\subsection{Midpoint and offset wavenumbers}

We can, if desired, convert source--receiver coordinates to midpoint--offset coordinates. Midpoint and offset are defined as that the midpoint 2-vector comprises the first two coordinates of the image function wavenumber \pmb{\textit{k}}\!:
\begin{equation}{\boldsymbol{k}}=\left(\tilde{\boldsymbol{k}}_m,\text{\textit{$k_z$}}\right)=\left(\text{\textit{$k_x$}},\text{\textit{$k_y$}},\text{\textit{$k_z$}}\right).\end{equation}

Equation (3.110) can be solved for frequency as a function of \(\tilde{\boldsymbol{k}}_m\), \(\tilde{\boldsymbol{k}}_h\), and \(k_z\):
\begin{equation}\omega =- \frac{\text{\textit{$k$}}_{\text{\textit{$z$}}}\text{\textit{$c$}}}{2}\sqrt{1+\text{\textit{$\tilde{\text{\textit{k}}}$}}_{\text{\textit{$m$}}}{}^2/\text{\textit{$k$}}_{\text{\textit{$z$}}}{}^2+\text{\textit{$\tilde{\text{\textit{k}}}$}}_{\text{\textit{$h$}}}{}^2/\text{\textit{$k$}}_{\text{\textit{$z$}}}{}^2+\left(\tilde{\boldsymbol{k}}_{\text{\textit{$m$}}}\pmb{\cdot} \tilde{\boldsymbol{k}}_{\text{\textit{$h$}}}\right){}^2/\text{\textit{$k$}}_{\text{\textit{$z$}}}{}^4}.
\end{equation}

Using this relation, the vertical wavenumbers \(k_{\text{\textit{gz}}}\) and \(k_{\text{\textit{sz}}}\) can also be expressed in terms of \(\tilde{\boldsymbol{k}}_m\), \(\tilde{\boldsymbol{k}}_h\), and \(k_z\):
\begin{equation}
k_gz=+ \frac{k_z}{2}\left(1-\frac{{\tilde{\boldsymbol k}}_m{\bf \cdot}{\tilde{\boldsymbol k}_h}}{k_z^2}\right),
\end{equation}
\begin{equation}
k_sz=- \frac{k_z}{2}\left(1+\frac{{\tilde{\boldsymbol k}}_m{\bf \cdot}{\tilde{\boldsymbol k}_h}}{k_z^2}\right).\end{equation}

The offset wavenumber can also be extended into a 3-vector:
\begin{equation}\pmb{\text{\textit{$\boldsymbol{k}$}}}_{\text{\textit{${h}$}}}=\text{\textit{$\pmb{\text{\textit{$\boldsymbol{k}$}}}_g$}}+\text{\textit{$\pmb{\text{\textit{$\boldsymbol{k}$}}}_s.$}}\end{equation}

The offset 3-vector is normal to \pmb{\textit{k}}\!, hence lies in the plane of specular reflectors. Its \textit{z}-component is
\begin{equation}k_hz=-\frac{\tilde{\boldsymbol k}_m\cdot
{\tilde{\boldsymbol k}_h}}{2k_z}.\end{equation}


\subsection{Image amplitude and phase}

We now turn attention to the amplitude-phase function $\mathcal{A}$. In the spirit of Section 3.2.2, we choose $\mathcal{A}$ to give us the ``right''
data function in the case of a single dipping reflector, and use the principle of superposition\index{superposition!principle of} to claim we have found the right answer for an arbitrary image function. We choose the \textit{x}-axis to align with reflector dip, so that
\begin{equation}\text{\textit{$R$}}({\mathbf{\mathit{x}}},\gamma ) =\text{\textit{$R$}}_0(\gamma )\delta \left(\left(\text{\textit{$z$}}-\text{\textit{$z$}}_0\right)\cos
 \alpha -\text{\textit{$x$}}\text{\textit{$ $}}\sin  \alpha \right).\end{equation}

The space--time formula for the reflection response \textit{D} can be determined from the method of images.  The desired response is a delta function
\begin{equation}\text{\textit{$D$}}\left(\tilde{\boldsymbol{x}}_g|\tilde{\boldsymbol{x}}_s;\text{\textit{$t$}}\right)\text{
 }= \frac{\text{\textit{$R$}}_0(\gamma ) \delta (\text{\textit{$t$}}-2\text{\textit{$r$}}/\text{\textit{$c$}})}{2\text{\textit{$r$}}},\end{equation}
where \textit{r} is half the total distance from source to reflector to receiver. The amplitude decay of $1{/}r$ in the impulse response is referred to as spherical divergence. This is an asymptotic formula in that we have inserted a plane-wave specular reflection coefficient
\(R_0(\gamma )\) into the formula. Consequently, we will be content with an asymptotic relation between (3.128) and (3.129), or its frequency-domain
equivalent
\begin{equation}\text{\textit{$D$}}\left(\tilde{\boldsymbol{x}}_g|\tilde{\boldsymbol{x}}_s;\omega
\right)\text{  }=e^{i 2\omega  \text{\textit{$r/c$}}} \frac{\text{\textit{$R$}}_0(\gamma )}{2\text{\textit{$r$}}}.\end{equation}
In terms of wavenumber, \textit{R} is
\begin{eqnarray}
\text{\textit{$R$}}(\pmb{\text{\textit{$k$}}},\gamma ) & = & \text{\textit{$R$}}_0(\gamma )\int \int \int d\text{\textit{$x$}}d\text{\textit{$y$}}d\text{\textit{$z$}}\text{\textit{$
$}}\delta \left(\left(\text{\textit{$z$}}-\text{\textit{$z$}}_0\right)\cos  \alpha -\text{\textit{$x$}}\text{\textit{$ $}}\sin  \alpha \right)e^{-i\left(\text{\textit{$k_xx$}}+\text{\textit{$k_yy$}}\text{\textit{$+$}}\text{\textit{$k_zz$}}\right)}
\nonumber\\
& = & (2\pi )^2\text{\textit{$R$}}_0(\gamma )e^{-i \text{\textit{$k_z$}}\text{\textit{$z$}}_0}\delta \left(\text{\textit{$k_y$}}\right)\delta \left(\text{\textit{$k_x$}}\cos
 \alpha +\text{\textit{$k_z$}}\sin  \alpha \right).
\end{eqnarray}
Invoking (3.103), the frequency--wavenumber reflection response \(D\) must be
\begin{equation}
\text{\textit{$D$}}\left(\tilde{\boldsymbol{k}}_{\text{\textit{$g$}}}|\tilde{\boldsymbol{k}}_{\text{\textit{$s$}}};\omega \right)\text{  }= \frac{ (2\pi )^2\text{  }\text{\textit{$R$}}_0(\gamma )}{\mathcal{A}\left(\tilde{\boldsymbol{k}}_{\text{\textit{$g$}}}|
\tilde{\boldsymbol{k}}_{\text{\textit{$s$}}};\omega
\right)}e^{-i \text{\textit{$k_z$}}\text{\textit{$z$}}_0}\delta \left(\text{\textit{$k_y$}}\right)\delta \left(\text{\textit{$k_x$}}\cos  \alpha
+\text{\textit{$k_z$}}\sin  \alpha \right)\!.\end{equation}
The delta functions in \textit{D} are satisfied when
\begin{eqnarray}\text{\textit{$k_y$}} & = & 0,\nonumber\\
\text{\textit{$k_x$}} & = & - \text{\textit{$k_{z }$}}\tan  \alpha ,\nonumber\\
1+\tilde{\boldsymbol{k}}_{\text{\textit{$m$}}}{}^2/
\text{\textit{$k$}}_{\text{\textit{$z$}}}{}^2 & = & \sec ^2\alpha ,\nonumber\\
\tan ^2\gamma & = & \frac{\text{\textit{$k_{\text{hx}}$}}{}^2+\text{\textit{$k_{\text{hy}}$}}{}^2\cos
^2\alpha }{\text{\textit{$k$}}_{\text{\textit{$z$}}}{}^2}.\end{eqnarray}
In space-frequency, the data function becomes
\begin{eqnarray}
\text{\textit{$D$}}\left(\tilde{\boldsymbol{x}}_g|\tilde{\boldsymbol{x}}_s; \omega
\right) & = & \frac{1}{(2\pi )^2}\int \int d^2\text{\textit{$k$}}_{\text{\textit{$g$}}}d^2\text{\textit{$k$}}_{\text{\textit{$s$}}} \frac{ \text{\textit{$R$}}_0(\gamma
)}{\mathcal{A}\left(\tilde{\boldsymbol{k}}_{\text{\textit{$g$}}}|
\tilde{\boldsymbol{k}}_{\text{\textit{$s$}}};\omega
\right)}e^{i \left(\tilde{\boldsymbol k}_g\cdot \tilde{\boldsymbol x}_g-{\tilde{\boldsymbol k}_s\cdot \tilde{\boldsymbol x}_s}-k_zz_0\right)}
\nonumber\\
 && \times  \delta \left(\text{\textit{$k_y$}}\right)\delta \left(\text{\textit{$k_x$}}\cos  \alpha +\text{\textit{$k_z$}}\sin  \alpha \right).
\end{eqnarray}
Make a change of integration variables to midpoint and offset wavenumber:
\begin{align}\text{\textit{$D$}}\left(\tilde{\boldsymbol{x}}_g|\tilde{\boldsymbol{x}}_s;\omega
\right)\text{  } & =  \frac{1}{4(2\pi )^2}\int \int d^2\text{\textit{$k$}}_{\text{\textit{$m$}}}d^2\text{\textit{$k$}}_{\text{\textit{$h$}}}
\text{  }\frac{\text{\textit{$R$}}_0(\gamma
)}{\mathcal{A}\left(\tilde{\boldsymbol{k}}_{\text{\textit{$g$}}}|
\tilde{\boldsymbol{k}}_{\text{\textit{$s$}}};\omega
\right)}e^{i (\tilde{\boldsymbol k}_g\cdot \tilde{\boldsymbol{x}}_g-\tilde{\boldsymbol k}_s\cdot \tilde{\boldsymbol{x}}_s-k_zz_0)}\nonumber\\
& \times  \delta \left(\text{\textit{$k_y$}}\right)
 \delta \left(\text{\textit{$k_x$}}\cos  \alpha
+\text{\textit{$k_z$}}\sin  \alpha \right)
\nonumber\\
& = \frac{1}{16\pi ^2}\int d^2\text{\textit{$k$}}_{\text{\textit{$h$}}} \frac{1}{\cos  \alpha } \frac{\text{\textit{$R$}}_0(\gamma )}{\mathcal{A}\left(\tilde{\boldsymbol{k}}_{\text{\textit{$g$}}}|
\tilde{\boldsymbol{k}}_{\text{\textit{$s$}}};\omega
\right)}e^{i (\tilde{\boldsymbol k}_h\cdot \tilde{\boldsymbol{x}}_h-k_zz_0+x_m\tan  \alpha )}.\end{align}

Given that \(k_y=0\) and \(k_x=-k_z\tan  \alpha \), equation (3.114) can be solved for \(k_z\) as a function of $\omega $ and \(\tilde{\boldsymbol{k}}_h\):
\begin{equation}\text{\textit{$k$}}_{\text{\textit{$z$}}}=- \text{Sign}(\omega ) \sqrt{\left(\frac{4\omega ^2}{\text{\textit{$c$}}^2}-\text{\textit{$k$}}_{\text{\textit{hy}}}{}^2\right)\cos
^2\alpha -\text{\textit{$k_{\text{hx}}$}}{}^2} .\end{equation}

In the far field, the remaining integrals can be evaluated via stationary phase.\index{stationary phase approximation} The result is
\begin{equation}\text{\textit{$D$}}\left(\tilde{\boldsymbol{x}}_g|\tilde{\boldsymbol{x}}_s;\omega
\right) \simeq - \frac{i \omega  \text{\textit{$z$}}_{\text{\textit{$m$}}}\cos  \alpha }{4\pi \text{\textit{$c$}}\text{\textit{$ $}}\text{\textit{$r
$}}^2}\text{  }\frac{\text{\textit{$R$}}_0(\underset{{}^{\wedge}}{\gamma })}{\mathcal{A}\left(
\text{\textit{$\underset{{}^{\wedge}}{\tilde{\boldsymbol{k}}}$}}{}_{\text{\textit{$g$}}}|
\text{\textit{$\underset{{}^{\wedge}}{\tilde{\boldsymbol{k}}}$}}{}_{\text{\textit{$s$}}};\omega
\right)}e^{2i \omega  \text{\textit{$r/c$}}},\end{equation}
with \(\underset{{}^{\wedge}}{\tilde{\boldsymbol{k}}}{}_g\) and \(\underset{{}^{\wedge}}{\tilde{\boldsymbol{k}}}{}_s\) the stationary
values of the source and receiver wavenumbers, and \(\underset{{}^{\wedge}}{\gamma }\) the stationary value of angle of incidence.

Comparing equation (3.137) to the dipping-layer formula (3.130), we discover that the amplitude-phase function must be
\begin{equation} \mathcal{A}\left(\underset{{}^{\wedge}}{\tilde{\boldsymbol{k}}}{}_g|\underset{{}^{\wedge}}{\tilde{\boldsymbol{k}}}{}_s;\omega \right) \equiv - \frac{i \omega
\text{\textit{$z$}}_{\text{\textit{$m$}}}\cos  \alpha }{2\pi  \text{\textit{$c$}}\text{\textit{$ $}}\text{\textit{$r$}}\text{\textit{$ $}}}.\end{equation}

The stationary value of angle of incidence is given by the relation
\begin{equation}\cos  \underset{{}^{\wedge}}{\gamma }=\frac{\text{\textit{$z$}}_{\text{\textit{$m$}}}\cos  \alpha }{\text{\textit{$r$}}},\end{equation}

The stationary values of the various wavenumbers are
\begin{equation}\text{\textit{$\underset{{}^{\wedge}}{k}{}_{\text{hx}}$}}=\frac{2\omega }{\text{\textit{$c$}}}\frac{\cos ^2\text{\textit{$ $}}\alpha
 \text{\textit{$x_h$}}}{\text{\textit{$r$}}},\quad\text{\textit{$\underset{{}^{\wedge}}{k}{}_{\text{hy}}$}}=\frac{2\omega }{\text{\textit{$c$}}}\frac{\text{\textit{$y_h$}}}{\text{\textit{$r$}}},\end{equation}
and
\begin{equation}\text{\textit{$\underset{{}^{\wedge}}{k}{}_z$}}=-\frac{2\omega }{\text{\textit{$c$}}}\frac{\cos ^2 \alpha  \text{\textit{$z_m$}}}{\text{\textit{$r$}}}.\end{equation}

With this last relation, we can write the amplitude-phase function (3.138) in terms of vertical wavenumber as
\begin{equation} \mathcal{A}\left(\underset{{}^{\wedge}}{\tilde{\boldsymbol{k}}}{}_g|
\underset{{}^{\wedge}}{\tilde{\boldsymbol{k}}}{}_s;\omega \right) =\frac{ i{\text{\textit{$\underset{{}^{\wedge}}{\text{\textit{$k$}}}$}}}{}_{\text{\textit{$z$}}}}{4\pi
 \cos  \alpha \text{\textit{$ $}}}.\end{equation}

The dip angle $\alpha $ is also expressible in terms of wavenumbers as
\begin{equation}\text{  }\cos  \alpha \text{\textit{$ $}}=\frac{1}{ \sqrt{1+\text{\textit{$\underset{{}^{\wedge}}{\tilde{k}}{}_m$}}{}^2/\text{\textit{$\underset{{}^{\wedge}}{k}{}_z$}}{}^2}},\end{equation}
which allows us to write the dip-independent expression
\begin{equation} \mathcal{A}\left(\tilde{\boldsymbol{k}}_{\text{\textit{$g$}}}|\tilde{\boldsymbol{k}}_{\text{\textit{$s$}}};\omega \right) =\frac{ i\text{\textit{$k$}}_{\text{\textit{$z$}}} \sqrt{1+\text{\textit{$\tilde{k}_m$}}{}^2/\text{\textit{$k_z$}}{}^2}}{4\pi
}.\end{equation}

With this choice for $\mathcal{A}$, the operator (3.103) will (asymptotically, anyway) yield a true-amplitude reflectivity\index{reflectivity} function for reflectors
of any dip or orientation, which is to say, it will convert any superposition of dipping-layer data functions of the form (3.129) into a superposition
of true-amplitude reflections.

\section{Prestack reverse-time migration}

One can implement prestack reverse-time migration\index{migration!prestack reverse time} by continuing reflection data backward in time while continuing forward in time a wave from the
source, forming two wavefields \(P_s(x,y,z,t)\) and \(P_r(x,y,z,t)\). We know that these two wavefields are related through the reflectivity\index{reflectivity},
and would like to infer from them the reflectivity function. Try the following, and assume that at any point (\textit{x, y, z}) in space,
\begin{equation}\text{\textit{$P_r(x,y,z,t)$}}\text{\textit{$ $}}=\text{\textit{$ $}}\text{\textit{$R$}}\text{\textit{$(x,y,z)$}} \text{\textit{$P_i$}}\text{\textit{$(x,y,z,t)$}}\text{\textit{$
$}}+\text{\textit{$n(x,y,z,t).$}}\end{equation}

That is, at every point in space and time, the reflected wave has a component proportional to reflectivity\index{reflectivity} times the incident wave at the same point
in space and time. The reflected wave may have additional components, but we hope they will not correlate with the incident wave at (\textit{
x, y, z, t}). This suggests a weighted integral over time:
\begin{eqnarray}
\int _{-\infty }^{\infty }d\text{\textit{$t$}} \text{\textit{$P_i$}}\text{\textit{$(x,y,z,t)$}}\text{\textit{$ $}}\text{\textit{$P_r$}}\text{\textit{$(x,y,z,t)$}}\text{\textit{$
$}}&=&\text{\textit{$ $}}\text{\textit{$R$}}\text{\textit{$(x,y,z)$}} \int _{-\infty }^{\infty }d\text{\textit{$t$}}\text{\textit{$ $}}\text{\textit{$P_i$}}\text{\textit{$(x,y,z,t)$}}^2\text{\textit{$
$}}\nonumber\\
&&+\int _{-\infty }^{\infty }d\text{\textit{$t$}} \text{\textit{$P_i$}}\text{\textit{$(x,y,z,t)$}}\text{\textit{$ $}}\text{\textit{$n$}}\text{\textit{$(x,y,z,t).$}}\qquad\qquad
\end{eqnarray}

In the hope that the last term is negligible, we have
\begin{equation}\text{\textit{$R(x,y,z)$}}\simeq \frac{\int _{-\infty }^{\infty }d\text{\textit{$t$}} \text{\textit{$P_i$}}\text{\textit{$(x,y,z,t)$}}\text{\textit{$
$}}\text{\textit{$P_r$}}\text{\textit{$(x,y,z,t)$}}\text{\textit{$ $}}}{\int _{-\infty }^{\infty }d\text{\textit{$t$}}\text{\textit{$ $}}\text{\textit{$P_i$}}\text{\textit{$(x,y,z,t)$}}^2\text{\textit{$
$}}}.\end{equation}

By neglecting the last term we open the door to artifacts. However, it is reasonable to hope for decent results.  To test the basic concept,
look at a simple case for which we can produce analytic results, namely, a constant-velocity layered medium. Put a single source at position \(\left(x_s,y_s,0\right)\),
and endow it with a Gaussian wavelet (Appendix E, Section E.3) of characteristic width \(t_w.\) The forward-continued response to this source is
\begin{equation}\text{\textit{$P_i(x,y,z,t)$}}\text{\textit{$ $}}= \frac{1}{\sqrt{\pi }\text{\textit{$t_w$}}\tau _i}e^{-\left(\text{\textit{$t$}}-\tau
_i\right){}^2/\text{\textit{$t_w$}}{}^2},\end{equation}
with
\begin{equation}\tau _i=\sqrt{\left(\text{\textit{$x$}}-\text{\textit{$x_s$}}\right){}^2+\left(\text{\textit{$y$}}-\text{\textit{$y_s$}}\right){}^2+\text{\textit{$z$}}^2}/\text{\textit{$c.$}}\end{equation}

With the given normalization, the integral over all time of \(P_i\) is \(1/r_i\). In the limit of small \(t_w\) it approaches the
delta function form
\begin{equation}\text{\textit{$P_i(x,y,z,t)$}}\text{\textit{$ $}}\underset{\text{\textit{$t_w$}}\to 0}{\to } \delta \left(\text{\textit{$t$}}-\tau
_i\right)/\text{\textit{$\tau _i$}}.\end{equation}

The forward-continued wave is oblivious to reflectors in the medium.

Suppose the medium does have flat reflectors at depths \(a_j\). If reflectivity\index{reflectivity} is small enough that transmission effects are negligible, then
the reflection response at zero depth is
\begin{equation}\text{\textit{$P_r$}}\text{\textit{$($}}\text{\textit{$x$}}\text{\textit{$,$}}\,\text{\textit{$y$}}\text{\textit{$,$}}\,0\text{\textit{$,$}}\,\text{\textit{$t$}}\text{\textit{$)$}}\text{\textit{$
$}}=\sum _{\text{\textit{$j$}}} \text{\textit{$ $}}\frac{\text{\textit{$R$}}_{\text{\textit{$j$}}}}{\sqrt{\pi }\text{\textit{$t_w$}}\tau _{\text{\textit{$j$}}}\text{\textit{$($}}\text{\textit{$x$}}\text{\textit{$,$}}\text{\textit{$y$}}\text{\textit{$,$}}0\text{\textit{$)$}}}e^{-\left(\text{\textit{$t$}}-\tau
_{\text{\textit{$j$}}}\text{\textit{$($}}\text{\textit{$x$}}\text{\textit{$,$}}\text{\textit{$y$}}\text{\textit{$,$}}0\text{\textit{$)$}}\right){}^2/\text{\textit{$t_w$}}{}^2},\end{equation}
where
\begin{equation}\tau _{\text{\textit{$j$}}}\text{\textit{$($}}\text{\textit{$x$}}\text{\textit{$,$}}\,\text{\textit{$y$}}\text{\textit{$,$}}\,0\text{\textit{$)$}}=\sqrt{\left(\text{\textit{$x$}}-\text{\textit{$x_s$}}\right){}^2+\left(\text{\textit{$y$}}-\text{\textit{$y_s$}}\right){}^2+4\text{\textit{$a$}}_{\text{\textit{$j$}}}{}^2}/\text{\textit{$c.$}}\end{equation}

If we backward-propagate the reflected wave, we obtain
\begin{equation}\text{\textit{$P_r(x,y,z,t)$}}\text{\textit{$ $}}=\text{\textit{$ $}}\sum _{\text{\textit{$j$}}} \text{\textit{$ $}}\frac{\text{\textit{$R$}}_{\text{\textit{$j$}}}}{\sqrt{\pi
}\text{\textit{$t_w$}}\tau _{\text{\textit{$j$}}}\text{\textit{$(x,y,z)$}}}e^{-\left(\text{\textit{$t$}}-\tau _{\text{\textit{$j$}}}\text{\textit{$(x,y,z)$}}\right){}^2/\text{\textit{$t_w$}}{}^2},\end{equation}
where
\begin{equation}\tau _{\text{\textit{$j$}}}\text{\textit{$(x,y,z)$}}=\sqrt{\left(\text{\textit{$x$}}-\text{\textit{$x_s$}}\right){}^2+\left(\text{\textit{$y$}}-\text{\textit{$y_s$}}\right){}^2+\left(2\text{\textit{$a$}}_{\text{\textit{$j$}}}\text{\textit{$-$}}\text{\textit{$z$}}\text{\textit{$)$}}\right.{}^2}/\text{\textit{$c.$}}\end{equation}

This expression holds both above and beneath the reflectors because the backward-propagated field, unlike the actual reflected wavefield, does not see the reflectors.

Now we try the reflectivity\index{reflectivity} formula (3.147). The denominator in this case is
\begin{eqnarray} \int _{-\infty }^{\infty }d\text{\textit{$t$}}\text{\textit{$ $}}\text{\textit{$P_i$}}\text{\textit{$(x,y,z,t)$}}^2\text{\textit{$
$}} & = & \frac{1}{\pi  \text{\textit{$t_w$}}{}^2\tau _{\text{\textit{$i$}}}{}^2}\int _{-\infty }^{\infty }d\text{\textit{$t$}}\text{\textit{$ $}}e^{-2\left(\text{\textit{$t$}}-\tau
_i\right){}^2/\text{\textit{$t_w$}}{}^2}\nonumber\\
& = & \frac{1}{\sqrt{2\pi } \text{\textit{$t_w$}}\tau _{\text{\textit{$i$}}}{}^2}.\end{eqnarray}

The numerator evaluates to
\begin{eqnarray}
\int _{-\infty }^{\infty }d\text{\textit{$t$}} \text{\textit{$P_i$}}\text{\textit{$(x,y,z,t)$}}\text{\textit{$ $}}\text{\textit{$P_r$}}\text{\textit{$(x,y,z,t)$}}\text{\textit{$
$}} & = & \sum _{\text{\textit{$j$}}} \frac{\text{\textit{$R$}}_{\text{\textit{$j$}}}}{\pi  \text{\textit{$t_w$}}{}^2\tau _{\text{\textit{$i$}}}\tau _{\text{\textit{$j$}}}}\int
_{-\infty }^{\infty }d\text{\textit{$t$}} e^{-\left[\left(\text{\textit{$t$}}-\text{\textit{$\tau _i$}}\right){}^2+\left(\text{\textit{$t$}}-\tau
_{\text{\textit{$j$}}}\right){}^2\right]/\text{\textit{$t_w$}}{}^2}
\nonumber\\
 & = & \frac{1}{\sqrt{2\pi } \text{\textit{$t_w$}}\tau _{\text{\textit{$i$}}}}\sum _{\text{\textit{$j$}}} \frac{\text{\textit{$R$}}_{\text{\textit{$j$}}}}{\tau
_{\text{\textit{$j$}}}} e^{\left.-\left(\tau _{\text{\textit{$i$}}}-\tau _{\text{\textit{$j$}}}\right){}^2\right/2 \text{\textit{$t_w{}^2$}}}.
\end{eqnarray}

Combining numerator and denominator,
\begin{equation}\text{\textit{$R(x,y,z)$}}\simeq \sum _{\text{\textit{$j$}}} \frac{\tau _{\text{\textit{$i$}}}}{\tau _{\text{\textit{$j$}}}}\text{\textit{$R$}}_{\text{\textit{$j$}}}
e^{\left.-\left(\tau _{\text{\textit{$i$}}}-\tau _{\text{\textit{$j$}}}\right){}^2\right/2 \text{\textit{$t_w{}^2$}}}.\end{equation}

Since \(t_w\) is small, the exponential functions in the reflectivity\index{reflectivity} estimate should be sharply peaked around points where \(\tau _i=\tau _j.\)
At the \(j{\text{th}}\) such point, the peak occurs where
\begin{equation}\sqrt{\left(\text{\textit{$x$}}-\text{\textit{$x_s$}}\right){}^2+\left(\text{\textit{$y$}}-\text{\textit{$y_s$}}\right){}^2+\text{\textit{$z$}}^2}=\sqrt{\left(\text{\textit{$x$}}-\text{\textit{$x_s$}}\right){}^2+\left(\text{\textit{$y$}}-\text{\textit{$y_s$}}\right){}^2+\left(2\text{\textit{$a$}}_{\text{\textit{$j$}}}\text{\textit{$-$}}\text{\textit{$z$}}\text{\textit{$)$}}\right.{}^2},\end{equation}
the obvious solution to which is \(z=a_j.\) For \textit{z} near \(a_j,\) define \(\epsilon =z-a_j,\) in terms of which
\begin{eqnarray}
\tau _{\text{\textit{$i$}}}-\tau _{\text{\textit{$j$}}} & = & \frac{1}{\text{\textit{$c$}}}\left(\sqrt{\left(\text{\textit{$x$}}-\text{\textit{$x_s$}}\right){}^2
+\left(\text{\textit{$y$}}-\text{\textit{$y_s$}}\right){}^2+\left(\text{\textit{$a_j$}}+\epsilon
\right){}^2}\right.\nonumber\\
 &&- \left.\sqrt{\left(\text{\textit{$x$}}-\text{\textit{$x_s$}}\right){}^2
+\left(\text{\textit{$y$}}-\text{\textit{$y_s$}}\right){}^2+\left(\text{\textit{$a_j$}}-{\epsilon}\right){}^2}\right)
\nonumber\\
& = & \frac{2 \text{\textit{$a_j$}}\epsilon }{\text{\textit{$c$}}\sqrt{\left(\text{\textit{$x$}}-\text{\textit{$x_s$}}\right){}^2+\left(\text{\textit{$y$}}-\text{\textit{$y_s$}}\right){}^2+\text{\textit{$a_j$}}{}^2}}+O\left(\epsilon
^3\right).
\end{eqnarray}

For \(t_w\) sufficiently small, only the term lowest order in $\epsilon $ is significant, whence (replacing the ratio \(\tau _i/\tau _j\) by 1),
\begin{equation}\text{\textit{$R(x,y,z)$}}\simeq \sum _{\text{\textit{$j$}}} \text{\textit{$R$}}_{\text{\textit{$j$}}}\,\text{exp}\left[-\frac{2\left(\text{\textit{$z-a_j$}}\text{\textit{$)$}}\right.{}^2}{\text{\textit{$c
$}}^2\text{\textit{$t_w$}}{}^2\left(1+\left(\left(\text{\textit{$x$}}-\text{\textit{$x_s$}}\right){}^2+\left(\text{\textit{$y$}}-\text{\textit{$y_s$}}\right){}^2\right)/\text{\textit{$a_j$}}{}^2\right)}\right].
\end{equation}

In this instance, the formula (3.147) gives a bandlimited reflectivity\index{reflectivity} estimate, with responses proportional to the reflection coefficients \(R_j\)
centered at the true reflector depths \(a_j\). The characteristic width of the migrated response increases with distance between source and observation
point, and decreases with reflector depth \(a_j\). Summing images from multiple source points would mitigate this effect, and could also help
with any artifacts.

Equation (3.147) made the assumption that reflectivity depends on spatial position only, and not on incoming and reflected\index{angle!reflected} angles. Unlike asymptotic
methods, in which waves know the direction they are traveling, a finite-difference reverse-time operator does not so easily distill directional information.
This is not to say that an amplitude-preserving reverse-time migration algorithm is impossible, but it would require a more sophisticated imaging
algorithm than (3.147).

Amplitude and artifact issues aside, the ability of reverse-time migration to handle two-way wave propagation through complex structure enhances
its appeal. A 3-D prestack reverse-time migration will be computationally very intensive and will consume large amounts of computer
memory and bandwidth. Even so, as computer power continues to increase, interest in reverse-time migration continues to~grow.


